% concatenated from arXiv_V3.tex
%%%%%%%%%%%%%%%%%%%%%%%%%%%%%%%%%%%% 
% 
%  F. FORSTNERIC  
% 
%  The Oka principle for sections of fibre bundles of 
% H\"older--Zygmund class on strongly pseudoconvex domains
% 
%  AMS-LaTeX 1.2 file for journals, based on amsart.cls. 
%
%  Submitted to Acta Math. Sinica on 24 May 2026.
%
%  Accepted for publication in Acta Math. Sinica, English Series, on 11 Sept. 2026
%
%  The final version sent on 14 Sept 2026.
%
%%%%%%%%%%%%%%%%%%%%%%%%%%%%%%%%%% 

\documentclass[11pt]{amsart}

\usepackage{amsfonts,amssymb,amscd,amstext,mathrsfs,mathtools}
\usepackage[utf8]{inputenc}
\usepackage{hyperref}
\usepackage{verbatim}

\usepackage{times}
\usepackage{enumerate}


\textwidth = 157mm
\textheight = 239mm
\evensidemargin=0mm
\oddsidemargin=0mm
\hoffset=4mm
\voffset=-25mm
\parskip =1mm
\parindent = 6mm
\linespread{1.09}
\pagestyle{plain}

%% To do not cut the words
\pretolerance=2000
\tolerance=3000

\numberwithin{equation}{section}



%%%%%%%%%% 
%%%%%%%%%% 
%%%%%%%%%% 
%%%%%%%%%% 
%%%%%%%%%% 
%%%%%%%%%% 

\newtheorem{theorem}{Theorem}[section] 
\newtheorem{proposition}[theorem]{Proposition} 
\newtheorem{claim}[theorem]{Claim} 
\newtheorem{lemma}[theorem]{Lemma} 
\newtheorem{corollary}[theorem]{Corollary} 
%\newtheorem{remark}[theorem]{Remark} 

\theoremstyle{definition} 
\newtheorem{definition}[theorem]{Definition} 
\newtheorem{remark}[theorem]{Remark} 
\newtheorem{question}[theorem]{Question} 
\newtheorem{problem}[theorem]{Problem} 
\newtheorem{example}[theorem]{Example} 
\newtheorem{conjecture}[theorem]{Conjecture} 


%
% BOLDFACE CAPITAL
%
\newcommand\bA{\mathbf{A}} 
\newcommand\bB{\mathbf{B}} 
\newcommand\bC{\mathbf{C}} 
\newcommand\bD{\mathbf{D}} 
\newcommand\bE{\mathbf{E}} 
\newcommand\bH{\mathbf{H}} 
\newcommand\bI{\mathbf{I}} 
\newcommand\bK{\mathbf{K}} 
\newcommand\bL{\mathbf{L}} 
\newcommand\bM{\mathbf{M}} 
\newcommand\bN{\mathbf{N}} 
\newcommand\bO{\mathbf{O}} 
\newcommand\bP{\mathbf{P}} 
\newcommand\bR{\mathbf{R}} 
\newcommand\bS{\mathbf{S}} 
\newcommand\bT{\mathbf{T}} 
\newcommand\bU{\mathbf{U}} 
\newcommand\bV{\boldface{V}} 


%%%%%%%%%% 
%%%%%%%%%% 
%%%%%%%%%% 
%%%%%%%%%% CALIGRAPHIC CAPITAL 
%%%%%%%%%% 
%%%%%%%%%% 
%%%%%%%%%% 

\newcommand\Acal{\mathcal{A}} 
\newcommand\Bcal{\mathcal{B}} 
\newcommand\Ccal{\mathcal{C}} 
\newcommand\Dcal{\mathcal{D}} 
\newcommand\Ecal{\mathcal{E}} 
\newcommand\Fcal{\mathcal{F}} 
\newcommand\Hcal{\mathcal{H}} 
\newcommand\Ical{\mathcal{I}} 
\newcommand\Kcal{\mathcal{K}} 
\newcommand\Lcal{\mathcal{L}} 
\newcommand\Mcal{\mathcal{M}} 
\newcommand\Ncal{\mathcal{N}} 
\newcommand\Ocal{\mathcal{O}} 
\newcommand\Pcal{\mathcal{P}} 
\newcommand\Rcal{\mathcal{R}} 
\newcommand\Scal{\mathcal{S}} 
\newcommand\Tcal{\mathcal{T}} 
\newcommand\Ucal{\mathcal{U}} 
\newcommand\Vcal{\mathcal{V}} 

%%%%%%%%%% 
%%%%%%%%%% 
%%%%%%%%%% SMALL BOLDFACE 
%%%%%%%%%% 
%%%%%%%%%% 
%%%%%%%%%% 

\newcommand\ba{\mathbf{a}} 
\newcommand\be{\mathbf{e}} 
\newcommand\bu{\mathbf{u}} 
\newcommand\bv{\mathbf{v}} 
\newcommand\bx{\mathbf{x}} 
\newcommand\by{\mathbf{y}} 
\newcommand\bw{\mathbf{w}} 
\newcommand\bz{\mathbf{z}} 

%%%%%%%%%% 
%%%%%%%%%% 
%%%%%%%%%% MATH SCRIPT 
%%%%%%%%%% 
%%%%%%%%%% 
%%%%%%%%%% 

\newcommand\Ascr{\mathscr{A}} 
\newcommand\Bscr{\mathscr{B}} 
\newcommand\Cscr{\mathscr{C}} 
\newcommand\Dscr{\mathscr{D}} 
\newcommand\Fscr{\mathscr{F}} 
\newcommand\Iscr{\mathscr{I}} 
\newcommand\Lscr{\mathscr{L}} 
\newcommand\Oscr{\mathscr{O}} 
\newcommand\Sscr{\mathscr{S}} 
\newcommand\Uscr{\mathscr{U}} 
\newcommand\Gscr{\mathscr{G}} 
\newcommand\Yscr{\mathscr{Y}} 
\newcommand\Zscr{\mathscr{Z}} 
\newcommand\Pscr{\mathscr{P}} 
\newcommand\Qscr{\mathscr{Q}} 
\newcommand\Vscr{\mathscr{V}} 

%%%%%%%%%% 
%%%%%%%%%% 
%%%%%%%%%% MATH BLACKBOARD 
%%%%%%%%%% 
%%%%%%%%%% 
%%%%%%%%%% 

\newcommand\A{\mathbb{A}} 
\newcommand\B{\mathbb{B}} 
\newcommand\C{\mathbb{C}} 
\newcommand\D{\overline{\mathbb D}} 
\newcommand\CP{\mathbb{CP}} 
\renewcommand\D{\mathbb D} 
\newcommand\G{\mathbb{G}} 
\renewcommand\H{\mathbb{H}} 
\newcommand\K{\mathbb{K}} 
\newcommand\M{\mathbb{M}} 
\newcommand\N{\mathbb{N}} 
\renewcommand\P{\mathbb{P}} 
\newcommand\R{\mathbb{R}} 
\newcommand\RP{\mathbb{RP}} 
\renewcommand\S{\mathbb{S}} 
\newcommand\T{\mathbb{T}} 
\newcommand\U{\mathbb{U}} 
\newcommand\Z{\mathbb{Z}} 



%%%%%%%%%% 
%%%%%%%%%% 
%%%%%%%%%% FRAKTUR 
%%%%%%%%%% 
%%%%%%%%%% 
%%%%%%%%%% 

\newcommand\ggot{\mathfrak{g}} 
\newcommand\hgot{\mathfrak{h}} 
\newcommand\igot{\mathfrak{i}} 
\newcommand\jgot{\mathfrak{j}} 
\newcommand\kgot{\mathfrak{k}} 
\renewcommand\igot{\mathfrak{i}} 
\newcommand\mgot{\mathfrak{m}} 
\newcommand\pgot{\mathfrak{p}} 
\newcommand\qgot{\mathfrak{q}} 
\newcommand\fgot{\mathfrak{f}} 

\newcommand\Agot{\mathfrak{A}} 
\newcommand\Igot{\mathfrak{I}} 
\newcommand\Lgot{\mathfrak{L}} 
\newcommand\Mgot{\mathfrak{M}} 
\newcommand\Ngot{\mathfrak{N}} 
\newcommand\Pgot{\mathfrak{P}} 
\newcommand\Ogot{\mathfrak{O}} 
\newcommand\Ygot{\mathfrak{Y}} 
\newcommand\Ggot{\mathfrak{G}} 
\newcommand\Hgot{\mathfrak{H}} 
\newcommand\Zgot{\mathfrak{Z}} 

% 
% typewriter 
% 

\newcommand\jtt{\mathtt{j}} 
\newcommand\ktt{\mathtt{k}} 
\newcommand\ttI{\mathtt{I}} 

% 
% e, i, zero - mathmode 
% 
\newcommand\E{\mathrm{e}} 
\newcommand\I{\imath} 
\renewcommand\imath{\igot} 
\newcommand\zero{\mathbf{0}} 

\newcommand\Psf{\mathsf{P}} 


% 
% arrows 
% 
\newcommand\hra{\hookrightarrow} 
\newcommand\lra{\longrightarrow} 
\newcommand\Lra{\Longrightarrow} 
\newcommand\longhookrightarrow{\ensuremath{\lhook\joinrel\relbar\joinrel\rightarrow}} 

% 
% tilde, widehat, di, dibar 
% 
\newcommand\wt{\widetilde} 
\newcommand\wh{\widehat} 
\newcommand\di{\partial} 
\newcommand\dibar{\overline\partial} 
\newcommand{\nad}[2]{\genfrac{}{}{0pt}{}{#1}{#2}} 


% 
% abbreviations 
% 

\newcommand\dist{\mathrm{dist}} 
\renewcommand\span{\mathrm{span}} 
\newcommand\length{\mathrm{length}} 
\newcommand\reg{\mathrm{reg}} 
\newcommand\Fr{\mathrm{Fr}} 
\newcommand\Id{\mathrm{Id}} 
\newcommand\pr{\mathrm{pr}} 
\newcommand\rank{\mathrm{rank}} 

\def\Ell1{\mathrm{Ell_1}} 
\def\CEll1{\mathrm{CEll_1}} 
\def\CAP{\rm{CAP}} 
\def\rmax{\mathrm{rmax}} 
\def\cB{\overline{\B}}
\def\Jst{J_{\mathrm{st}}} 

\newcommand\Flux{\mathrm{Flux}} 
\newcommand\GCMI{\mathrm{GCMI}} 
\newcommand\CMI{\mathrm{CMI}} 
\newcommand\NC{\mathrm{NC}} 

\newcommand\Psh{\mathrm{Psh}} 
\newcommand\MPsh{\mathrm{\mathfrak{M}Psh}} 
\newcommand\MSPsh{\mathfrak{MS}\mathrm{Psh}} 
\newcommand\NPsh{\mathrm{\mathfrak{N}Psh}} 
\newcommand\NSPsh{\mathfrak{NS}\mathrm{Psh}} 
\newcommand\MS{\mathfrak{MS}} 

\newcommand\Co{\mathrm{Co}} 
\newcommand\Conv{\mathrm{Conv}} 
\newcommand\supp{\mathrm{supp}} 
\newcommand\Qcal{\mathcal{Q}} 
\newcommand\tr{\mathrm{tr}} 
\newcommand\Hess{\mathrm{Hess}} 
\newcommand\Aut{\mathrm{Aut}} 
\newcommand\Aff{\mathrm{Aff}} 
\newcommand\Pic{\mathrm{Pic}} 
\newcommand\Hom{\mathrm{Hom}}
\newcommand\Lin{\mathrm{Lin}}
\newcommand\Lip{\mathrm{Lip}}

\newcommand\Ocalc{\overline{\mathcal{O}}}
\newcommand\Ocalcl{\Ocalc_{\mathrm{loc}}}
\newcommand\Pcalc{\overline{\mathcal{P}}}
\newcommand\Rcalc{\overline{\mathcal{R}}}

\newcommand\Oscrc{\overline{\mathscr{O}}}
\newcommand\Oscrcl{\Oscrc_{\mathrm{loc}}}

\newcommand\boldA{\mathbf A}

%\renewcommand\i{\imath} 

% The following makes 2020 to appear in the MSC
\makeatletter
\@namedef{subjclassname@2020}{\textup{2020} Mathematics Subject Classification}
\makeatother


%%%%%%%%%% 
%%%%%%%%%% 
%%%%%%%%%% 
%%%%%%%%%% 
%%%%%%%%%% 
%%%%%%%%%% 

\begin{document} 

\title{The Oka principle for holomorphic fibre bundles of 
H\"older--Zygmund classes on strongly pseudoconvex domains}

\author{Franc Forstneri{\v c}}

\address{Franc Forstneri{\v c}, Faculty of Mathematics and Physics, University of Ljubljana, Jadranska 19, 1000 Ljubljana, Slovenia}

\address{Franc Forstneri{\v c}, Institute of Mathematics, Physics, and Mechanics, Jadranska 19, 1000 Ljubljana, Slovenia}

\email{franc.forstneric@fmf.uni-lj.si}

\subjclass[2020]{Primary 32Q56; secondary 32L05, 32T15, 46J15}

% 32E10 Stein spaces
% 32G05 Deformations of complex structures
% 32H02 Holomorphic mappings, (holomorphic) embeddings and related questions in several complex variables 
% 32L05 Holomorphic bundles and generalizations
% 32L20 Vanishing theorems  
% 32Q28 Stein manifolds
% 32Q56 Oka principle and Oka manifolds
% 32T15: Strongly pseudoconvex domains
% 32V05(2000?now) CR structures, CR operators, and generalizations
% 32V20(2000?now) Analysis on CR manifolds
% 32W05 ?��� and ?���-Neumann operators 
% 46J15 Banach algebras of differentiable or analytic functions

\date{24 May 2026. This version: 14 September 2026}

\keywords{fibre bundle, Oka manifold, Oka principle, 
H\"older--Zygmund space} 

\begin{abstract} 
Let \(\overline \Omega\) be a compact strongly pseudoconvex domain
with smooth boundary in a Stein manifold, 
and let \(h:Z\to \overline \Omega\) be a 
fibre bundle of H\"older--Zygmund class \(\Lambda^r\), \(r>0\), 
which is holomorphic over \(\Omega\).
Assuming that the fibre is an Oka manifold, we prove that 
every continuous section \(f_0:\overline \Omega\to Z\) 
is homotopic to a section \(f_1:\overline \Omega\to Z\) 
of class \(\Lambda^r(\overline \Omega)\) which is holomorphic on \(\Omega\).
We also establish the parametric $h$-principle in this context.  
As an application, we obtain the Oka principle for the classification
of vector bundles and principal bundles of H\"older--Zygmund classes
on such domains. 
\end{abstract}

\maketitle

\centerline{\em In Memory of the 100th Birthday of Professor Lu Qikeng}

\setcounter{tocdepth}{1}
\tableofcontents

%
%  SECTION: Introduction
%
\section{Introduction}
\label{sec:intro}

A complex manifold $Y$ is said to be an {\em Oka manifold} 
(see \cite{Forstneric2009CR} and \cite[Chap.\ 5]{Forstneric2017E}) if maps
$X\to Y$ from any Stein manifold $X$ satisfy the h-principle, also called
the Oka principle. This means in particular that any continuous map 
$f_0:X\to Y$ is homotopic to a holomorphic map $f_1:X\to Y$;
if $f_0$ is holomorphic on a neighbourhood of a compact holomorphically
convex subset $K$ of $X$ and on a closed complex subvariety 
$X'$ of $X$, then a homotopy $f_t:X\to Y$ $(t\in [0,1])$ 
from $f_0$ to a holomorphic map $f_1$ can be chosen to 
be fixed on $X'$ and to consist of maps which are holomorphic on a 
neighbourhood of $K$ and uniformly close to $f_0$ on $K$. 
The analogous results hold
for sections of holomorphic fibre bundles with Oka fibres, and for sections
of elliptic holomorphic submersions $Z\to X$ onto a Stein space.
For the theory of Oka manifolds and Oka maps, see 
\cite{Forstneric2023Indag,Forstneric2026ICM,Forstneric2017E,ForstnericLarusson2011}. 
Classical examples of Oka manifolds include complex homogeneous 
manifolds (see Grauert \cite{Grauert1957I,Grauert1957II} and 
\cite[Proposition 5.6.1]{Forstneric2017E}) and, more generally, 
Gromov elliptic manifolds 
(see \cite{Gromov1989} and \cite[Corollary 5.6.14]{Forstneric2017E}).
Oka manifolds $Y$ are characterised by the approximation
property for holomorphic maps $K\to Y$ from (neighbourhoods of)
compact convex sets $K\subset\C^n$ by entire maps
$\C^n\to Y$ (the {\em convex approximation property}, CAP);
see \cite{Forstneric2009CR} and \cite[Theorem 5.4.4]{Forstneric2017E}), 
and by {\em convex relative ellipticity}, CRE  
(see Kusakabe \cite[Theorem 1.3]{Kusakabe2021IUMJ} and
\cite[Definition 1.5 and Theorem 1.6]{Forstneric2026ICM}).
It has recently been shown that every projective Oka manifold is 
elliptic \cite{ForstnericLarusson2025MRL}, but there exist 
noncompact Oka manifolds which fail to be elliptic 
\cite{Kusakabe2024AM,Forstneric2023Indag}. 
Modern Oka theory has diverse applications.

In the present paper, we establish the following 1-parametric
Oka principle for sections of fibre bundles of 
H\"older--Zygmund classes $\Lambda^r_\Oscr(\overline\Omega)$ 
with Oka fibres on compact strongly pseudoconvex domains 
$\overline \Omega$ with Stein interior.
The definition of the Banach algebra 
$\Lambda^r_\Oscr(\overline\Omega)$ for any real number $r>0$ 
is given in Subsect.\ \ref{ss:HZspaces}, manifold-valued maps
of H\"older--Zygmund classes are introduced in Subsect.\ 
\ref{ss:manifold-valued}, and fibre bundles of H\"older--Zygmund classes
are introduced in Subsect.\ \ref{ss:Lr-bundles}.

%
%  MAIN THEOREM
%
\begin{theorem}\label{th:main}
Assume that $\Omega$ is a relatively compact, strongly pseudoconvex 
domain with smooth boundary $b\Omega$ in a Stein manifold $X$, $r>0$, 
and $h:Z\to \overline \Omega$ is a fibre bundle of H\"older--Zygmund 
class $\Lambda^r(\overline \Omega)$ which is holomorphic on $\Omega$.
Assuming that the fibre of $h$ is an Oka manifold, every continuous
section $f_0\in \Gamma(\overline \Omega,Z)$ of $h$ is homotopic to a section 
$f\in \Gamma^r_\Oscr(\overline \Omega,Z)$ of class 
$\Lambda^r(\overline \Omega)$ which is holomorphic on $\Omega$. 
Furthermore, every homotopy of continuous sections 
$\{f_t\}_{t\in [0,1]}\in \Gamma(\overline \Omega,Z)$ with 
$f_0, f_1 \in \Gamma^r_\Oscr(\overline \Omega,Z)$ can be deformed with
fixed ends to a homotopy in $\Gamma^r_\Oscr(\overline \Omega,Z)$. 
\end{theorem}

Theorem \ref{th:main} is proved in Sect.\ \ref{sec:proofmain}, where we
also provide a fully parametric version, Theorem \ref{th:mainpar}.
The proof also gives an approximation result common in Oka 
theory: if the section $f_0$ in Theorem \ref{th:main} is holomorphic
on a neighbourhood of a compact holomorphically convex subset
$K\subset \Omega$, then the homotopy from $f_0$ to a holomorphic
section $f_1$ can be chosen to consist of sections that are holomorphic
on a neighbourhood of $K$ and approximate $f_0$ uniformly on $K$.
A stronger approximation theorem, with $K\subset\overline \Omega$
holomorphically convex in $\overline \Omega$, is also possible;
see \cite[Theorem 6.1]{DrinovecForstneric2008FM} where the
analogue of this result is established for fibre bundles 
$Z\to \overline \Omega$ of class 
\[
	\Ascr^r(\overline\Omega)
	=\{f\in \Cscr^r(\overline \Omega):f|_\Omega \in\Oscr(\Omega)\}
\]
with $r\in \Z_+=\{0,1,2,\ldots\}$. Here, $\Oscr(\Omega)$ denotes the space
of holomorphic functions on $\Omega$.
A motivation for the present generalisation of the mentioned results
from \cite{DrinovecForstneric2008FM}  is that the spaces $\Lambda^r$ 
behave better than the $\Cscr^k$ or Lipschitz spaces for many  
operators considered in analysis, as was already noticed by Zygmund in 1945. 
I was asked by Andrei Teleman (private communication, January 2026) 
whether Theorem \ref{th:main} holds.
After the completion of this work, he obtained a result
on extending principal holomorphic fibre bundles of 
H\"older--Zygmund classes \cite{Teleman2026}.

Theorem \ref{th:main} and the other results of the 
paper also hold for any compact complex manifold 
$\overline \Omega$ with Stein interior and smooth 
strongly pseudoconvex boundary. Indeed, by 
Ohsawa \cite{Ohsawa1984LNM}, Heunemann \cite{Heunemann1986III}, 
and Catlin \cite{Catlin1988}, such a manifold 
holomorphically embeds as a smoothly bounded strongly 
pseudoconvex domain in a Stein manifold.

We refer to Subsect.\ \ref{ss:HZspaces} for the definition
and basic properties of H\"older--Zygmund spaces
$\Lambda^r(\overline \Omega)$ on a bounded Lipschitz domain
$\Omega$ in a smooth Riemannian manifold. 
For $r=k+\alpha$ with $k\in\Z_+$ and $0<\alpha<1$, 
the space $\Lambda^r(\overline \Omega)$ coincides with the  
H\"older space $\Cscr^{k,\alpha}(\overline \Omega)$. On the other hand,
for integer values $r=k+1\in \{1,2,\ldots\}$ the Zygmund space
$\Lambda^{k+1}(\overline \Omega)$ properly contains 
the Lipschitz space $\Cscr^{k,1}(\overline \Omega)$ 
of functions whose derivatives up to order $k$ 
are Lipschitz continuous on $\overline \Omega$
(see \cite[Example 4.3.1, p.\ 148]{SteinEM1970}). 
The difference between the Lipschitz and the Zygmund
spaces is that, in the definition of the Zygmund 
$\Lambda^1$ norm of a function $f$, one replaces the first 
difference $\Delta_h f(x)$ (see \eqref{eq:Delta})
by the second difference $\Delta^2_h f(x)$ (see \eqref{eq:Delta2}).
The Zygmund class $\Lambda^{1}$
is a natural substitute for the Lipschitz class $\Cscr^{0,1}=\Lip^1$ 
in many contexts. There are continuous strict embeddings among the 
classical and the H\"older--Zygmund scales on any Lipschitz domain;
see \eqref{eq:scale}. 

An important fact used in the proof of Theorem \ref{th:main} is that the 
canonical (Kohn) solution operator for the $\dibar$-equation
is bounded on $\Lambda^r(\overline \Omega)$
when $\Omega$ is a smoothly bounded strongly pseudoconvex 
domain; see Beals et al.\ \cite{BealsGreinerStanton1987}.
The precise result that we shall use is stated as Theorem \ref{th:dibar}.

On the way to Theorem \ref{th:main}, 
we obtain several other results of independent interest.
Given a Lipschitz domain $\Omega\Subset X$ in a complex manifold $X$
and a real number $r>0$, let
\[
	\Lambda^r_\Oscr(\overline \Omega) = 
	\{f\in \Lambda^r(\overline \Omega): f|_\Omega\in \Oscr(\Omega)\}.
\] 
Similarly we define the mapping spaces $\Oscr(\Omega,Y)$ and 
$\Lambda^r_\Oscr(\overline \Omega,Y) \subset 
\Lambda^r(\overline \Omega,Y)$ for any complex manifold $Y$;
see Subsect.\ \ref{ss:manifold-valued}.  
Given a compact set $K\subset X$, 
$\Oscr(K)$ denotes the space of restrictions to $K$ 
of holomorphic functions in open neighbourhoods of $K$, endowed
with the inverse limit topology. The notation
$\Oscr(K,Y)$ is used for the space of maps of this kind
to a complex manifold $Y$. We have the following approximation theorem. 

%
%	APPROXIMATION
%
\begin{theorem}\label{th:approximation}
Let $\Omega$ be a relatively compact, smoothly bounded, strongly 
pseudoconvex domain in a Stein manifold $X$, and let 
$Y$ be a complex manifold. Every map 
in $\Lambda^r_\Oscr(\overline \Omega,Y)$, $r>0$, 
can be approximated in the $\Lambda^r(\overline \Omega,Y)$ 
topology by maps in $\Oscr(\overline\Omega,Y)$. 
\end{theorem}

Theorem \ref{th:approximation} is proved in Sect.\ \ref{sec:approximation};
see also the parametric version in Theorem \ref{th:approximation-par}. 
The analogue of Theorem \ref{th:approximation} is known for spaces 
$\Ascr^r(\overline\Omega)$ with $r\in\Z_+$; 
see \cite[Theorem 2.9.2,\ p.\ 87]{HenkinLeiterer1984} 
and \cite[Theorem 24, p.\ 165]{FornaessForstnericWold2020}.  
For approximation of manifold-valued maps in 
$\Ascr^r(\overline\Omega,Y)$, where $Y$ is a complex manifold
and $r\in\Z_+$,  
by holomorphic maps from open neighbourhoods of $\overline\Omega$
in $X$, see \cite[Theorem 1.2]{DrinovecForstneric2008FM},
\cite[Theorem 8.11.4]{Forstneric2017E}, and 
\cite[Corollary 9, p.\ 178]{FornaessForstnericWold2020}.
The last mentioned result in \cite{FornaessForstnericWold2020}
is a more general Mergelyan-type approximation theorem 
on strongly admissible sets in Stein manifolds.

We also have the following result generalising 
\cite[Theorem 1.1 (i)]{Forstneric2007AJM}.
The proof in \cite[Sect.\ 2]{Forstneric2007AJM} also applies to spaces
$\Lambda^r_\Oscr(\overline \Omega,Y)$.

\begin{theorem}\label{th:manifold}
Assume that $\Omega$ is 
a relatively compact, smoothly bounded, strongly 
pseudoconvex domain in a Stein manifold and 
$Y$ is a complex manifold. For every $r>0$ 
the space $\Lambda^r_\Oscr(\overline \Omega,Y)$ 
is a complex Banach manifold.
The tangent space $T_f \Lambda^r_\Oscr(\overline \Omega,Y)$ 
at a point $f\in \Lambda^r_\Oscr(\overline \Omega,Y)$
is the space of sections of class $\Lambda^r_\Oscr(\overline \Omega)$
of the complex vector bundle $f^*TY\to \overline \Omega$.
\end{theorem}

In the proof of Theorem \ref{th:main}, we shall need several
results concerning vector bundles $E\to\overline\Omega$ of
class $\Lambda^r_\Oscr(\overline\Omega)$ 
on smoothly bounded strongly pseudoconvex Stein domains
$\Omega$; see Sect.\ \ref{sec:vectorbundles}. 
In particular, we obtain Theorem A for such bundles; 
see Theorems \ref{th:TheoremA} and \ref{th:TheoremApar}. 
The proof is based on the Cartan splitting lemma 
for maps of H\"older--Zygmund classes to a complex Lie group;
see Lemmas \ref{lem:Cartan}, \ref{lem:Cartan-par}
and Remark \ref{rem:Liegroup}. 
The main technical results used in the proof of  Theorem \ref{th:main}
are a splitting lemma (see Lemma \ref{lem:splitting}) 
and a gluing lemma (see Lemma \ref{lem:gluing})
for sprays of sections of class $\Lambda^r_\Oscr$. 

An application of Theorem \ref{th:main} is 
the following Oka principle for vector bundles of class 
$\Lambda^r_\Oscr(\overline\Omega)$, proved in Sect.\ \ref{sec:OPVB}. 
The analogous result for principal fibre bundles 
is given by Theorem \ref{th:OPprincipalbundles}. 

\begin{theorem}\label{th:OPVB}
Let $\Omega$ be a relatively compact, smoothly bounded,
strongly pseudoconvex domain in a Stein manifold.
The following hold for every real number $r>0$.
\begin{enumerate}[\rm (i)]
\item Every topological complex vector bundle on $\overline\Omega$
is isomorphic to a vector bundle of class 
$\Lambda^r_\Oscr(\overline\Omega)$.
\item If a pair of vector bundles 
of class $\Lambda^r_\Oscr(\overline\Omega)$ are isomorphic as
topological complex vector bundles, then they are also
isomorphic as $\Lambda^r_\Oscr(\overline\Omega)$ vector bundles.
\end{enumerate}
\end{theorem}

This result is classical for vector bundles on Stein spaces;
see Grauert \cite{Grauert1958MA}, with the special case 
of line bundles due to Oka \cite{Oka1939}. 
(See also Leiterer \cite{Leiterer1990} and 
\cite[Theorem 5.3.1]{Forstneric2017E}.) For vector bundles
of classes $\Ascr^k(\overline\Omega)$, $k\in \Z_+$,  
see Leiterer \cite{Leiterer1976I,Leiterer1976II}
and Heunemann \cite[Theorem 2]{Heunemann1986I}.


%
%  OTHER MAPPING SPACES
%
\begin{remark}\label{rem:Fcal}
The results of this paper apply to a wider class of mapping spaces.
Assume that $\Fcal$ is a contravariant functor from the category 
of compact smooth manifolds $M$ with boundary 
to the category of Banach algebras of $\C$-valued functions 
on them, satisfying the following conditions:
%
\begin{enumerate}[\rm (a)]
\item 
$\Cscr^\infty(M) \subset \Fcal(M)\subset \Cscr(M)$ and  
both inclusions are continuous. 
\item 
A smooth map $\Phi:M\to M'$ induces a homomorphism 
$\Phi^*:\Fcal(M') \to \Fcal (M)$ of Banach algebras by
$f\mapsto f\circ \Phi$ for $f\in \Fcal(M')$.
\item 
Postcomposition by a smooth function $\C\to\C$ induces
a continuous selfmap of $\Fcal(M)$.
\end{enumerate}
%
These properties imply that the topology on $\Fcal(M)$ 
can be defined via local charts, and hence the
definition of these classes extends to differential forms and other tensor 
fields on $M$. Furthermore, we can introduce vector bundles
and more general fibre bundles of class $\Fcal$; 
see Palais \cite{Palais1968,Palais1971} and the 
references in \cite{Forstneric2007AJM}.
%
When $M=\overline \Omega$ is a compact complex manifold with 
smoothly boundary, we define
\[
	\Fcal_\Oscr(\overline \Omega)=
	\{f\in \Fcal(\overline \Omega): f|_\Omega\in\Oscr(\Omega)\}.
\]
The analogous definition yields the mapping space
$\Fcal_\Oscr(\overline\Omega,Y)$ for any complex manifold $Y$.
Conditions (a)--(c) on the functor $\Fcal$ clearly 
imply the following properties:
%
\begin{itemize} 
\item 
A smooth map $\Phi:\overline\Omega \to \overline\Omega'$
which is holomorphic on $\Omega$ induces a homomorphism 
$\Phi^*:\Fcal_\Oscr(\overline\Omega') \to \Fcal (\overline\Omega)$ by
$f\mapsto f\circ \Phi$, $f\in \Fcal_\Oscr(\overline\Omega')$.
%
\item 
Postcomposition by an entire function $g:\C\to\C$ induces
a continuous selfmap of $\Fcal_\Oscr(\overline \Omega)$.
\end{itemize}
%
To conditions (a)--(c) we add the following condition:
%
\begin{enumerate}[\rm (d)]
\item There is a bounded linear operator
$T:\Fcal_{0,1}(\overline\Omega) \to \Fcal(\overline\Omega)$
satisfying $\dibar \, T(\alpha)=\alpha$ for any $(0,1)$-form
$\alpha\in \Fcal_{0,1}(\overline\Omega)$ with $\dibar\alpha=0$.
\end{enumerate}
%
In the proofs of our results, the operator $\dibar$ 
is only applied to functions
of the form $\chi f$ with $\chi\in \Cscr^\infty(\overline\Omega)$
and $f\in \Fcal_\Oscr(\overline\Omega)$. In this case, the derivative
$\dibar(\chi f)=f\dibar\chi \in \Fcal_{0,1}(\overline\Omega)$ 
is of the same class as $f$.

Inspections of proofs of our results show that 
they hold on Banach algebras $\Fcal_\Oscr(\overline\Omega)$,
and for vector and fibre bundles of this class, when
$\Fcal$ satisfies condition (a)--(d). 
Examples include H\"older spaces $\Cscr^{k,\alpha}(\overline\Omega)$ 
$(k\in\Z_+,\ 0<\alpha<1)$,
H\"older--Zygmund spaces $\Lambda^r(\overline\Omega)$ $(r>0)$, 
and Sobolev spaces $W^{k,p}(\overline \Omega)$
$(k\in \N,\ 1\le p<\infty,\ kp>\dim_\R \Omega)$, among others. 
\end{remark}


\begin{remark}[Added in the final revision]
The main results of the paper are stated for strongly pseudoconvex 
domains $\Omega$ with smooth boundaries, the reason being 
that Theorem \ref{th:dibar} requires smooth boundaries. 
However, the existence of a linear $\overline\partial$-solution operator 
gaining $1/2$ derivative (or with no loss of derivatives) suffices for our
construction. In the case of Sobolev spaces, a suitable result
under the weaker condition that $b\Omega$ is of class $\Cscr^2$ was 
obtained recently by Shi and Yao \cite{ShiYao2025}.
It is likely that a similar result holds for the Zygmund spaces.
\end{remark}

%
%
% 	SECTION: PRELIMINARIES
%
%
\section{Preliminaries on H\"older--Zygmund spaces}\label{sec:prelim}

In the first subsection, we recall the definition and basic properties of 
H\"older spaces $\Cscr^{k,\alpha}$ $(k\in\Z_+,\ 0<\alpha\le 1)$  
and H\"older--Zygmund spaces $\Lambda^r$ for any real $r>0$. 
In the second subsection, we explain how to define
manifold-valued maps of H\"older--Zygmund classes.
In the third subsection, we introduce vector bundles and more
general fibre bundles of these classes. In the fourth subsection,
we recall the results on the solution to the $\dibar$-equation in 
the spaces $\Lambda^r$, which will be used in the paper. 
For more information, we refer to the papers by 
Gong \cite[Sect.\ 5]{Gong2025TAMS}, Wallin \cite{Wallin1988}, 
and the monographs by 
Gilbarg and Trudinger \cite[Sect. 4.1]{GilbargTrudinger1983},
E.\ M.\ Stein \cite[Sect.\ V.4]{SteinEM1970}, 
and M.\ E.\ Taylor \cite{TaylorME2023}. 



%
%	HOLDER-ZYGMUND SPACES
%
\subsection{H\"older--Zygmund spaces} \label{ss:HZspaces}

Let $\Omega$ be a domain in $\R^n$. 
Given a function $f:\Omega\to\C$ and a vector 
$h\in \R^n\setminus\{0\}$, the first and the second difference of $f$ 
at $x\in \Omega$, with step $h$, are defined by 
\begin{eqnarray}\label{eq:Delta}
	\Delta_h f(x) &=& f(x+h)-f(x), 
	\quad \ \ \quad\qquad\qquad x+h\in \Omega; \\
	\label{eq:Delta2}
	\Delta^2_h f(x) &=& f(x+h)+f(x-h)-2f(x),\quad x+h,\ x-h\in \Omega.
\end{eqnarray}
Note that $\Delta^2_h f(x) =\Delta_h f(x) - \Delta_h f(x-h)$.
(Some authors use the definition 
$\Delta^2_h f(x)=\Delta_h \circ\Delta_h f(x) = f(x+2h)+f(x)-2f(x+h)$, 
which leads to the same function spaces.)
Given $x\in\Omega$ and $0<\alpha\le 1$, the function $f$ 
belongs to $\Cscr^{0,\alpha}(\Omega,x)$ if 
\[
	[f]_{\alpha,x} := 
	\sup_{h\ne 0,\ x+h\in \Omega} |h|^{-\alpha}|\Delta_h f(x)|
	= \sup_{y\in\Omega\setminus\{x\}} \frac{|f(y)-f(x)|}{|y-x|^\alpha}  
	< \infty.
\]
Such $f$ is said to be H\"older class $\alpha$ at the point $x$. 
For $\alpha=1$, $\Cscr^{0,\alpha}(\Omega,x)=\Lip^1(\Omega,x)$ is the 
Lipschitz space at $x$. The H\"older-$\alpha$ space on $\Omega$ is
\begin{equation}\label{eq:Calpha}
	\Cscr^{0,\alpha}(\Omega) =
	\bigl\{f:\Omega\to \C: \|f\|_{\Cscr^{(0,\alpha)}(\Omega)} = 
	\sup_{x\in \Omega} |f(x)| + \sup_{x\in \Omega}\, [f]_{\alpha,x}<\infty
	\bigr\}.
\end{equation}
For $\alpha=1$ we have the Lipschitz space
$\Cscr^{0,1}(\Omega)=\mathrm{Lip}^1(\Omega)$. 

Let $x=(x_1,\ldots,x_n)$ denote the coordinates on $\R^n$.
Given $\beta=(\beta_1,\ldots,\beta_n)\in\Z_+^n$, set 
$|\beta|=\beta_1+\cdots+\beta_n$ and  
$D^\beta f=\di^{|\beta|}f/\di x_1^{\beta_1} \cdots \di x_n^{\beta_n}$. 
For $k\in \Z_+$ we denote by $\Cscr^k(\Omega)$ the space
of $k$-times continuously differentiable functions $f$ on $\Omega$ with 
\[
	\|f\|_{\Cscr^k(\Omega)} 
	= \sum_{|\beta|\le k} \sup_{x\in \Omega} |D^\beta f(x)| < \infty.
\]
The H\"older space $\Cscr^{k,\alpha}(\Omega)$ for $k\in \Z_+$ and 
$0<\alpha\le 1$ is defined by 
\begin{eqnarray}\label{eq:Ckalpha}
	\Cscr^{k,\alpha}(\Omega) &=& 
	\Big\{f\in \Cscr^k(\Omega): D^\beta f\in \Cscr^{0,\alpha}(\Omega)
	\ \ \text{for all $\beta\in \Z_+^n$ with $|\beta|\le k$}, \\ 
	&& \quad 
	\label{eq:Ckalpha-norm}
	\|f\|_{\Cscr^{k,\alpha}(\Omega)} 
	:= \|f\|_{\Cscr^{k}(\Omega)} 
	+ \sum_{|\beta|=k} \|D^\beta f\|_{\Cscr^{0,\alpha}(\Omega)} <\infty 
	\Big\}.
\end{eqnarray}
For $\alpha=1$ one also writes 
$\Cscr^{k,1}(\Omega)= \mathrm{Lip}^{k,1}(\Omega)$.

Assume now that $\Omega\Subset \R^n$ is a bounded 
Lipschitz domain. This means that its boundary $b\Omega$
is locally at each point a Lipschitz graph over an 
affine hyperplane in $\R^n$, with the domain $\Omega$ lying on one side
of the graph. A function $f:\overline \Omega\to\C$ 
belongs to $\Cscr^k(\overline \Omega)$ for some $k\in\Z_+$ if it 
is the restriction to $\overline \Omega$ of a function 
$\tilde f\in \Cscr^k(\R^n)$. (For $k=0$, this coincides with the usual
definition of continuous functions by Tietze's extension theorem.) 
Given $\beta\in \Z_+^n$ with $|\beta|\le k$, we denote by 
$D^\beta f$ the restriction of $D^\beta \tilde f$ to $\overline \Omega$;
note that $D^\beta f$ is independent of the choice of the extension $\tilde f$
if $\Omega$ is a Lipschitz domain. 
Given $0<\alpha\le 1$, the H\"older space 
$\Cscr^{k,\alpha}(\overline \Omega)$ on the closed domain 
$\overline\Omega$ is defined by
\[
	\Cscr^{k,\alpha}(\overline \Omega)=
	\big\{f\in \Cscr^k(\overline \Omega):
	D^\beta f \in \Cscr^{0,\alpha}(\overline \Omega)\ \ 
	\text{for all $\beta\in \Z_+^n$ with $|\beta|\le k$} \big\},
\]
endowed with the norm $\|f\|_{\Cscr^{k,\alpha}(\Omega)}$ 
\eqref{eq:Ckalpha-norm}. It follows from definitions 
that we get the same norm by taking the suprema over
$x\in \overline\Omega$, so 
$\|f\|_{\Cscr^{k,\alpha}(\Omega)}=\|f\|_{\Cscr^{k,\alpha}(\overline\Omega)}$.

%
%  EXTENSION LEMMA
%
\begin{proposition}
\label{prop:extension}
If $\Omega$ is a bounded Lipschitz domain in $\R^n$
then every function in $\Cscr^{k,\alpha}(\Omega)$
$(k\in \Z_+,\ 0<\alpha\le 1)$ extends to a unique function  
$\tilde f\in \Cscr^{k,\alpha}(\overline \Omega)$.
\end{proposition}

\begin{proof}
For $k=0$ this holds by McShane's extension theorem  
\cite[Corollary 1, p.\ 840]{McShane1934}. If $k\ge 1$
and $f\in \Cscr^{k,\alpha}(\Omega)$, its partial derivatives
$D^\beta f$ of order $|\beta|=k$ belong to $\Cscr^{0,\alpha}(\Omega)$,
so they extend to functions $g_\beta=\Cscr^{0,\alpha}(\overline \Omega)$
by McShane's theorem. By Whitney's theorem \cite{Whitney1934TAMS}
(see also Malgrange \cite[Theorem 3.2 and Comp.\ 3.5]{Malgrange1966}), 
it follows that $f$ extends to a function $\tilde f\in \Cscr^k(\R^n)$
satisfying $D^\beta \tilde f= g_\beta$ on $\overline \Omega$ for 
all $|\beta|=k$. Thus, 
$\tilde f|_{\overline\Omega}\in \Cscr^{k,\alpha}(\overline\Omega)$.
Clearly, the extension $\tilde f|_{\overline\Omega}$ of $f$ is unique.
\end{proof}


We now recall the definition of H\"older--Zygmund spaces 
$\Lambda^r(\Omega)$ for any real number $r>0$. 
Write $r=k+\alpha$ with $k\in\Z_+$ and 
$0<\alpha\le 1$. If $r$ is not an integer (that is, $\alpha\ne 1$), set 
$\Lambda^r(\Omega)=\Cscr^{k,\alpha}(\Omega)$ and 
$\Lambda^r(\overline \Omega)=\Cscr^{k,\alpha}(\overline \Omega)$;
both spaces are endowed with the norm 
$\|f\|_{\Lambda^{k+\alpha}(\Omega)}:=\|f\|_{\Cscr^{k,\alpha}(\Omega)}$
\eqref{eq:Ckalpha-norm}. Assume now that $r\in \{1,2,\ldots\}$. 
Recall that the second difference $\Delta_h^2 f$ is given 
by \eqref{eq:Delta2}. We define
\begin{eqnarray}\label{eq:Lambda-1}
	\Lambda^1(\Omega) &=& \big\{f\in \Cscr(\Omega):
	\|f\|_{\Lambda^1(\Omega)} = 
	\sup_{x\in\Omega}|f(x)| + \sup_{x,x\pm h\in\Omega,\, h\ne 0}
	|h|^{-1} |\Delta_h^2 f(x)| <\infty \big\}, \\
	\label{eq:Lambda-r}
	\Lambda^r(\Omega) &=&
	\big\{f\in \Cscr^{r-1}(\Omega):
	\|f\|_{\Lambda^r(\Omega)} = \|f\|_{\Cscr^{r-1}(\Omega)} +
	\sum_{|\beta|=r-1} \|D^\beta f\|_{\Lambda^1(\Omega)} <\infty\big\},
\end{eqnarray}
where $r>1$ in \eqref{eq:Lambda-r}. 
(For $\Omega=\R^n$, see 
\cite[Propositions 8 and 9, pp.\ 146--147]{SteinEM1970}.)
The space $\Lambda^1(\R^n)$ is the classical 
Zygmund space on $\R^n$; omitting the term 
$\|f\|_\infty=\sup_{x\in\Omega}|f(x)|$ in \eqref{eq:Lambda-1}
gives the homogeneous Zygmund space.
Clearly, $\Lambda^1(\R^n)$ contains the Lipschitz space 
$\Cscr^{0,1}(\R^n) = \mathrm{Lip}^1(\R^n)$
but is not equal to it as shown by
\cite[Example 4.3.1, p.\ 148]{SteinEM1970}.
The norms \eqref{eq:Lambda-1}--\eqref{eq:Lambda-r}
are also called the {\em interior Zygmund norms} on $\Omega$.

For $r\in \N$, the Zygmund space $\Lambda^r(\overline \Omega)$ 
on a closed Lipschitz domain $\overline\Omega\subset\R^n$ 
cannot be defined in general by considering only values of 
functions in the points of $\overline \Omega$. 
Indeed, if $\Omega$ is strictly convex 
and $x,x+h\in b\Omega$ with $h\ne 0$, then 
$x-h\notin \overline \Omega$, so the second difference $\Delta^2_h f(x)$ 
is not defined. One possible definition is the following; see 
Wallin \cite[Definition 8, p.\ 105]{Wallin1988}. It can be used 
for any $r>0$. 
\begin{eqnarray}\label{eq:Lambda-r-Omega}
	\Lambda^r(\overline \Omega) &=& 
	\big\{f=\tilde f|_{\overline \Omega}: \tilde f\in  \Lambda^r(\R^n)\big\}, \\
	\|f\|_{\Lambda^r(\overline\Omega)} &=& 
	\label{eq:outer}
	\inf\big\{ \|\tilde f\|_{\Lambda^r(\R^n)}: \tilde f \in \Lambda^r(\R^n),\
	 \tilde f|_{\overline \Omega} =f\big\}. 
\end{eqnarray}
The norms \eqref{eq:outer} are called {\em exterior Zygmund norms} 
on $\Omega$. By a classical theorem due to Besov, Ilin and Nikolski 
\cite[Theorem 18.5, p.\ 63]{BesovIlinNikolskii1979},
every function in $\Lambda^r(\Omega)$ extends to a function in 
$\Lambda^r(\R^n)$.
Furthermore, the interior and the exterior definition of the spaces 
$\Lambda^r(\overline\Omega)$ and their norms are equivalent. 
Indeed, E.\ M.\ Stein \cite[Sect.\ VI.2]{SteinEM1970}
constructed a linear extension operator 
$E:\Cscr(\overline\Omega) \to \Cscr_0(\R^n)$ 
to the space of continuous functions on $\R^n$ with compact support 
such that $E:\Lambda^r(\overline\Omega)\to \Lambda^r(\R^n)$ 
is bounded for every $r>0$. Stein proved that his extension operator 
is bounded on Sobolev spaces; for H\"older--Zygmund spaces 
this was shown by Gong; see \cite[Corollary 5.3]{Gong2025TAMS}.
Furthermore, Shi and Yao \cite[Theorem 1.1]{ShiYao2024} 
proved that for $r=k+\alpha$ with $k\in \N$ and $0<\alpha\le 1$,
$\|f\|_{\Lambda^r(\Omega)}$ is equivalent to
$\sum_{j\le k} \|\nabla^j f\|_{\Lambda^\alpha(\Omega)}$.
Their proof uses Rychkov's universal extension 
for Besov spaces $B^r_{p,q}(\Omega)$; see \cite{Rychkov1999}
and note that $\Lambda^r(\Omega)=B^r_{\infty,\infty}(\Omega)$. 
For noninteger values $r>0$ and domains $\Omega$ 
with smooth boundaries, see also 
\cite[Lemma 6.37]{GilbargTrudinger1983} whose 
proof is based on Seeley's extension theorem \cite{Seeley1964}. 	

\begin{remark}
In \eqref{eq:Lambda-r-Omega}--\eqref{eq:outer},
one can also use extensions of $f$ to any domain $D\subset \R^n$
containing $\overline\Omega$. It is not clear whether these norms 
decrease to the interior norm $\|f\|_{\Lambda^r(\Omega)}$ 
\eqref{eq:Lambda-r} as $D$ shrinks to $\overline\Omega$.
\end{remark}
	
If $X$ is a smooth manifold and $\Omega$ is a Lipschitz domain in $X$ 
with compact closure, one defines the spaces $\Lambda^{r}(\Omega)$ 
and $\Lambda^{r}(\overline \Omega)$, $r>0$,
by using a finite system of smooth coordinate 
charts on $X$ covering $\overline \Omega$. 
The norms obtained in this way are comparable to one another;
see (3.11) and (3.12) in \cite[Lemma 3.1]{GanGong2024}.
(See also \cite{Palais1968,Palais1971} 
and the discussion and references in 
\cite{BourdaudCristoforis2002}, \cite[Sect.\ 2]{Forstneric2007AJM}, 
and \cite[Subsect.\ 3.2]{GongShi2024}.)
%
The space $\Lambda^r(\overline\Omega)$, $r>0$, 
is a commutative unital Banach algebra 
under pointwise multiplication, with Moser-type estimates for products, 
compositions, and inverses; see \cite[Lemmas 3.1--3.3]{Gong2019MA},
\cite[Lemma 6.3]{Gong2025TAMS}, and \cite{BourdaudCristoforis2002}. 
There are continuous strict embeddings among the 
classical and the H\"older--Zygmund scales on
any relatively compact Lipschitz domain $\Omega\Subset X$
in a smooth manifold $X$:
\begin{equation}\label{eq:scale}
	\Cscr^{k+1}\subset \Cscr^{k,1} \subset \Lambda^{k+1}\subset
	\Lambda^{k+\alpha} \subset 
	\Lambda^{k+\beta}\subset \Cscr^k,\quad 
	k\in\Z_+,\ 0<\beta<\alpha<1.
\end{equation}

\begin{definition}\label{def:local}
A function $f:X\to \C$ on a smooth manifold $X$ 
is said to be of {\em local class $\Lambda^r$}, $r>0$,  
if every point $x\in X$ has a compact Lipschitz neighbourhood
$K \subset X$ such that $f|_K \in \Lambda^r(K)$.
\end{definition}

The space $\Lambda^{r,loc}(X)$ of all functions of local class $\Lambda^r$ 
is a Fr\'echet space, and a Banach space if $X$ is compact 
(possibly with boundary). In the latter case, 
$\Lambda^{r,loc}(X)=\Lambda^{r}(X)$. 

%
%  Manifold valued maps of Zygmund classes
%
\subsection{Manifold-valued maps of Zygmund classes}
\label{ss:manifold-valued} 

Let $\Omega\Subset X$ be a Lipschitz domain in a smooth
manifold $X$, and let $Y$ be another smooth manifold.
We wish to define the space 
$
	\Lambda^r(\overline \Omega,Y)
$
of maps $f:\overline\Omega\to Y$ of Zygmund class 
$\Lambda^r(\overline\Omega, Y)$.
For this, we need to know that postcompositions of Zygmund 
class maps between Euclidean spaces by smooth maps 
preserve these classes. More precisely, if 
$u=(u_1,\ldots,u_m):\overline \Omega\to \R^m$, 
with $u_i\in \Lambda^r(\overline\Omega,\R^n)$ for $i=1,\ldots,m$, 
and $F$ is a smooth function on a neighbourhood of the range 
of $u(\overline \Omega)$ in $\R^m$, we must know that 
$F\circ u \in \Lambda^r(\overline \Omega)$.
For noninteger $r$ (i.e., for the H\"older spaces), this is 
classical; see e.g.\  \cite{GilbargTrudinger1983}.
When $r$ is an integer and $n=1$ (i.e., for scalar valued $u$), 
a suitable reference is \cite{BourdaudCristoforis2002}. 
A precise reference for the general vector-valued case 
with integer $r$ seems difficult to find.
For this reason, we include the following lemma.

%
%  Postcompositions by smooth functions
%
\begin{lemma}\label{lem:composition}
Let $\Omega\Subset\R^n$ and $D\Subset \R^m$ be Lipschitz domains, 
$u=(u_1,\ldots,u_m): \overline\Omega\to \bar D$ a map of class 
$\Lambda^r(\overline \Omega)$ for some $r\in \N$, and 
$F: \bar D \to \mathbb R$ a function of class $\Cscr^{r+1}(\overline D)$. 
Then, $F\circ u\in \Lambda^r(\overline \Omega)$.
\end{lemma}

\begin{proof}
Using the Stein extension operator 
\cite[Sect.\ VI.2]{SteinEM1970} and the results 
mentioned above, we may assume $F\in \Cscr^{r+1}(\R^m)$
and $u\in \Lambda^r(\R^n)$ both have compact support.
By $\|u\|_{\Lambda^r}$ we denote the Zygmund norm 
of the extended function, and likewise for $F$.

We first prove the lemma for $r=1$. Given points $x,y\in\R^n$,
consider the function 
\[
 	g(t)=F\big((1-t)u(x)+tu(x+y)\big) +F\big((1-t)u(x)+tu(x-y)\big).
\]
We have 
\[
	g(1)-g(0) = F\big(u(x+h)\big)+ F\big(u(x-y)\big)-2F(x) 
	= \Delta^2_y (F\circ u)(x).
\]
(See \eqref{eq:Delta2}.) There is a $t\in (0,1)$ such that $g(1)-g(0) =g'(t)$. We have
\begin{eqnarray*}
   g'(t) &=& \nabla F\big((1-t)u(x)+tu(x+y)\big) \big(u(x+y)-u(x)\big)\\
   && +\nabla F\big((1-t)u(x)+tu(x-y)\big) \big(u(x-y)-u(x)\big) \\
   &=& \nabla F\big((1-t)u(x)+tu(x+y)\big)\Delta^2_y u(x) \\
   && + \big[\nabla F\big((1-t)u(x)+tu(x-y)\big) 
   	- \nabla F\big((1-t)u(x)+tu(x+y)\big)\big] 
   				\big(u(x-y)-u(x)\big). 
\end{eqnarray*}
The first term after the last equality sign is clearly bounded by 
$\|F\|_{\Cscr^1}\|u\|_{\Lambda^1}|y|$. In the second term,
\begin{eqnarray*}
	&& \big|\nabla F((1-t)u(x)+tu(x-y)) - \nabla F((1-t)u(x)+tu(x+y))\big| \le \\
	&& \quad \le const \|F\|_{\Cscr^2} |u(x+y)-u(x-y)| 
	\le const \|F\|_{\Cscr^2} \|u\|_{\Lambda^\alpha}|y|^\alpha
\end{eqnarray*}
for any $\alpha \in (0,1)$, where the constant only depends on 
$\alpha$ and $m$. (We used that $\Lambda^1\subset \Lambda^\alpha$.
Without loss of generality we may take $\alpha=1/2$.)
Similarly, we estimate the last factor in the second term by 
\[
	|u(x-y)-u(x)| \le \|u\|_{\Lambda^{1-\alpha}} |y|^{1-\alpha}.
\]
Hence, the second term in the expression for $g'(t)$ is bounded by 
\[
	const \|F\|_{\Cscr^2} \|u\|_{\Lambda^\alpha}|y|^\alpha
	 \|u\|_{\Lambda^{1-\alpha}}|y|^{1-\alpha}
	 = const \|F\|_{\Cscr^2} \|u\|_{\Lambda^\alpha}
	 \|u\|_{\Lambda^{1-\alpha}} |y|.
\]
For $y \ne 0$ it follows that 
\[ 
	|y|^{-1} |\Delta^2_y (F\circ u)(x)| \le  
	\|F\|_{\Cscr^1}\|u\|_{\Lambda^1} +
	const \|F\|_{\Cscr^2} \|u\|_{\Lambda^\alpha}
	\|u\|_{\Lambda^{1-\alpha}},
\]
and hence $F\circ u\in \Lambda^1$. This proves the result for $r=1$. 
Assume now that $r>1$ and the lemma holds for functions in $\Lambda^{r-1}$.
Let $F\in \Cscr^{r+1}$ and $u=(u_1,\ldots,u_n)\in \Lambda^{r}$.
Writing the coordinates on $\R^n$ and $\R^m$ 
by $x=(x_1,\ldots,x_n)$ and $y=(y_1,\ldots,y_m)$, respectively, we have
\[
	\frac{\di}{\di x_i} F(u(x)) = 
	\sum_{j=1}^m \frac{\di F}{\di y_j}(u(x)) \frac{\di u_j}{\di x_i}(x),
	\quad i=1,\ldots,n.
\]
Since $\frac{\di F}{\di y_j}\in \Cscr^{r}$ for 
$j=1,\ldots,m$, the inductive assumption implies 
$\frac{\di F}{\di y_j}\circ u \in \Lambda^{r-1}$.
Furthermore, $\frac{\di u_j}{\di x_i}\in \Lambda^{r-1}$,
so $\frac{\di}{\di x_i} (F\circ u) \in \Lambda^{r-1}$.
It follows that $F\circ u\in \Lambda^{r}$ and the induction may continue.
\end{proof}


 \begin{remark}\label{rem:composition}
 With suitable adjustments, the proof of Lemma \ref{lem:composition}
 also applies to classes $\Lambda^r$ with noninteger $r$,
 that is, to H\"older classes. By elaborating on the proof, one can 
 obtain explicit estimates and see in particular that the postcomposition 
 by a smooth function on $\R^m$ defines a smooth operator 
 $\Lambda^r(\overline \Omega,\R^m)\to \Lambda^r(\overline \Omega)$. 
Lemma \ref{lem:composition} also lets us define the space 
 $\Lambda^{r,loc}(X,Y)$ of maps of local Zygmund class $\Lambda^r$
 between any pair of smooth manifolds.
 \end{remark}

%
%  Mapping spaces and fibre bundles
%
\subsection{Mapping spaces and fibre bundles of class 
$\Lambda^r_\Oscr(\overline\Omega)$} \label{ss:Lr-bundles}

On a smooth manifold $X$, one defines fibre bundles of class
$\Lambda^r$ by asking that the transition maps are
of local class $\Lambda^r$; see Definition \ref{def:local}
and Remark \ref{rem:composition}.
We now describe such bundles in more detail in 
the case of main interest for this paper. 

Assume that $X$ and $Y$ are complex manifolds and 
$\Omega\Subset X$ is an open Lipschitz domain. Set
\[
	\Lambda^r_\Oscr(\overline\Omega,Y)=
	\{f\in \Lambda^r(\overline\Omega,Y): 
	f\ \text{is holomorphic on}\ \Omega\},\quad r>0.
\]
In particular, 
$\Lambda^r_\Oscr(\overline\Omega,\C) =\Lambda^r_\Oscr(\overline\Omega)$.

Let $Y$ be a complex manifold. 
A holomorphic fibre bundle $h:Z\to \overline\Omega$ 
of class $\Lambda^r_\Oscr(\overline \Omega)$ with fibre $Y$  
has total space of the form 
\[
	Z= \bigsqcup_{i=1}^m U_i \times Y/\!\!\sim, 
\]  
where $\{U_i\}_{i=1}^m$ is an open covering of 
$\overline \Omega$ by relatively open Lipschitz domains  
$U_i\subset \overline \Omega$ (which can be chosen of
the form $U_i=\wt U_i\cap \overline \Omega$ with 
$\wt U_i$ an open Lipschitz domain in $\R^n$)
such that the intersections 
$U_{i,j}=U_i\cap U_j$ are also Lipschitz, and a point 
$(x,y)\in U_j \times Y$ $(x\in U_{i,j})$ 
is identified by the equivalence relation $\sim$ with the point
$(x,\phi_{i,j}(x,y))\in U_i\times Y$, where the map 
\[
	U_{i,j} \times Y \ni (x,y) \mapsto \phi_{i,j}(x,y)\in Y
\] 
is of class $\Lambda^{r,loc}_\Oscr(U_{i,j})$ in $x\in U_{i,j}$,
holomorphic in $y$, $\phi_{i,j}(x,\cdotp)\in \Aut(Y)$ 
for every $x\in U_{i,j}$, and the maps $\phi_{i,j}$ satisfy the usual
1-cocycle condition. 
A section $f:\overline \Omega\to Z$ of $h:Z\to \overline\Omega$ 
is given by a collection of maps $f_i:U_i\to Y$ satisfying the 
compatibility conditions
\[
	f_i(x) = \phi_{i,j}(x, f_j(x)),\quad x\in U_{i,j},\ i,j=1,\ldots,m.
\]
The section $f:\overline \Omega\to Z$ is of class 
$\Lambda^r_\Oscr(\overline \Omega)$ if and
only if $f_i\in \Lambda^r_\Oscr(U_i,Y)$ for $i=1,\ldots,m$. 
We denote by $\Gamma^r_\Oscr(\overline \Omega,Z)$ the space 
of sections of class $\Lambda^r_\Oscr(\overline \Omega)$
of the bundle $h:Z\to\overline\Omega$.
If $Y$ is a smooth manifold, the analogous definition gives
fibre bundles $Z\to \overline\Omega$ of class $\Lambda^r(\overline\Omega)$.

In particular, denoting by $GL_n(\C)$ the general linear group 
of rank $n$ over $\C$, a vector bundle $\pi:E\to\overline \Omega$
of class $\Lambda^r_\Oscr(\overline \Omega)$ and rank $n$
has fibre $\C^n$ and transition maps 
\[
	\phi_{i,j}(x,v)=A_{i,j}(x)v, \quad v\in\C^n,
\]
where the maps $A_{i,j}\in \Lambda^{r,loc}_\Oscr(U_{i,j},GL_n(\C))$
satisfy the 1-cocycle conditions
\[
	A_{i,i}=I,\quad A_{i,j}A_{j,i}=I,\quad 
	A_{i,j}A_{j,k}A_{k,i}=I\quad \text{for all}\ i,j,k=1,\ldots,m.
\]
Here, $I\in GL_n(\C)$ denotes the identity matrix.
By shrinking the sets $U_i$ in the covering, we can represent the same
vector bundle by transition maps 
$A_{i,j}\in \Lambda^{r}_\Oscr(U_{i,j},GL_n(\C))$.
A section of $E$ of class $\Gamma^r_\Oscr(\overline\Omega)$ 
is given by a collection of maps $f_i\in \Gamma^r_\Oscr(U_i,\C^n)$ 
satisfying the compatibility conditions
\[
	f_i(x) = A_{i,j}(x) f_j(x),\quad x\in U_{i,j},\ i,j=1,\ldots,m.
\]
Similarly one defines vector bundle morphisms of class 
$\Lambda^r_\Oscr(\overline \Omega)$, 
subbundles, quotient bundles, etc. 
We refer to Leiterer \cite{Leiterer1990} for a similar treatment of
vector bundles of class $\Ascr^r(\overline \Omega)$ 
for $r\in \Z_+$, i.e., of class $\Cscr^r(\overline \Omega)$ and
holomorphic on $\Omega$.


%
%  dibar solving operators on Lipschitz spaces
%
\subsection{Regularity of the canonical solution operator for the
$\dibar$-equation in H\"older--Zygmund spaces.}

Let $p\ge 0$ and $q\ge 1$ be integers, and
let $\Omega$ be a domain in a complex manifold $X$. 
The $\dibar$-problem on $\Omega$ asks for the existence 
and regularity properties of solutions of the equation 
$\dibar u=f$ for a differential $(p,q)$-form $f$ on $\Omega$
(or on its closure $\overline\Omega$) 
satisfying the necessary condition $\dibar f=0$. One of the
most successful techniques in this field is the $\dibar$--Neumann
method introduced in the pioneering works of Kohn \cite{Kohn1963,Kohn1964}; 
see the books by Folland and Kohn \cite{FollandKohn1972} and 
Chen and Shaw \cite{ChenShaw2001}, among others.

Let $\overline\Omega$ be a compact complex 
Hermitian manifold of dimension $n$ with $\Cscr^\infty$
boundary $b\Omega$ and Stein interior $\Omega$
such that at each point of $b\Omega$ the Levi form 
has at least $n-q$ positive eigenvalues. Let $N_q$ denote the 
Neumann operator for the complex Laplacian
\[
	\Box_q = \dibar^*_q\dibar_{q} +  \dibar_{q-1} \dibar^*_{q-1},
	\quad q\ge 1
\]
acting on $(p,q)$-forms which satisfy the $\dibar$-Neumann boundary conditions on $b\Omega$. This means that 
$N_q\Box_q = \Box_q N_q$ is the $L^2$ orthogonal projection onto 
the range of $\Box_q$. Then, $T_q= \dibar^*_{q-1}N_q$ is 
Kohn's canonical solution operator for the $\dibar$-equation 
$\dibar_{q-1}u = f$ with $f$ a $(p,q)$-form with $\dibar_q f=0$.
Denote by $\Lambda^r_{p,q}(\overline \Omega)$ the space
of $(p,q)$-forms on $\overline\Omega$ of class 
$\Lambda^r(\overline \Omega)$ for $r>0$. 
(Note that $(p,q)$-forms are sections of the vector bundle 
$\Lambda^{p,q}TX|_{\overline \Omega}$, so the spaces 
$\Lambda^r_{p,q}(\overline \Omega)$ are naturally defined.)
The following is a special case of 
\cite[Theorem 2]{BealsGreinerStanton1987} 
due to Beals, Greiner and Stanton.
See also Gong \cite{Gong2025TAMS}.

\begin{theorem}\label{th:dibar}
Let $\Omega$ be as above.
The canonical solution $T_q=\dibar_{q-1}^*N_q$ to the $\dibar$-equation 
$\dibar u=f$ for $f\in \Lambda^r_{p,q}(\overline \Omega)$ satisfying 
$\dibar f=0$ maps $\Lambda^r_{p,q}(\overline \Omega)\cap\ker \dibar$
boundedly to $\Lambda^{r+1/2}_{p,q-1}(\overline \Omega)$ 
for every $p\ge 0$, $q\ge 1$, and $r>0$. 
\end{theorem}

We shall use this result for $p=0$.
The gain of regularity for $1/2$ implies that the operator 
$T_q: \Lambda^r_{p,q}(\overline \Omega)\cap\ker \dibar
\to \Lambda^{r}_{p,q-1}(\overline \Omega)$ is compact, 
a fact which is useful in many applications. 

%
%
%	SECTION: APPROXIMATION OF MAPS PF CLASS \LAMBDA^r
%
%
\section{Approximation of maps of class 
$\Lambda^r_\Oscr$ on strongly pseudoconvex domains}
\label{sec:approximation}

In this section, we prove Theorem \ref{th:approximation} and its
parametric version, Theorem \ref{th:approximation-par}. 
We shall use the following notion of a Cartan pair; 
see \cite[Definitions 5.7.1 and 5.10.2]{Forstneric2017E}.

%
% CARTAN PAIR
%
\begin{definition}\label{def:Cartan-pair}
A pair $(A,B)$ of compact sets in a complex manifold $X$ is a 
{\em Cartan pair} if 
%
\begin{enumerate}[\rm (i)]
\item $A$, $B$, $D=A\cup B$ and $C=A\cap B$ are closures of 
smoothly bounded, strongly pseudoconvex domains with Stein
interior, and
%
\item 
$\overline{A\setminus B}\cap \overline{B\setminus A} =\varnothing$.
\end{enumerate}
A Cartan pair $(A,B)$ is {\em special} if there is a coordinate
neighbourhood of $B$ in $X$ in which the sets $B$ and $C$ 
are strongly convex. In this case, $B$ is said to be a convex bump 
attached to $A$.
\end{definition}

A more general notion of a Cartan pair in
\cite[Definition 5.7.1 (I)]{Forstneric2017E} is suitable 
when considering holomorphic maps on neighbourhoods of 
the respective sets, which may be any Stein compacts. 
We shall not need it in this paper since our 
focus is on mapping spaces on smoothly bounded domains. 

\begin{proof}[Proof of Theorem \ref{th:approximation}]
Since the domain $\Omega\Subset X$ is smoothly bounded
and strongly pseudoconvex, there is a smooth strongly plurisubharmonic
function $\rho$ on a neighbourhood $U\subset X$ of $\overline\Omega$
such that $\Omega=\{x\in U:\rho(x)<0\}$ and $d\rho_x\ne 0$
for every point $x\in b\Omega =\{\rho=0\}$. Pick $c>0$ such that
$\rho$ has no critical values on the interval $[0,c]$.
Set $A=\overline\Omega=\{\rho\le 0\}$ and $A'=\{\rho\le c\}$.
Given an open cover $\Ucal=\{U_j\}$ of $\overline{A'\setminus A}$ 
consisting of holomorphic coordinate charts $U_j\subset X$,  
\cite[Lemma 5.10.3]{Forstneric2017E} gives an increasing 
sequence of compact, smoothly bounded, 
strongly pseudoconvex domains 
\begin{equation}\label{eq:sequenceA}
	A_0:=A \subset A_1 \subset \cdots \subset A_m=A'
\end{equation}
for some $m\in\N$ such that for every $k=0,1,\ldots,m-1$ we have 
$A_{k+1}=A_k\cup B_k$, where $(A_k,B_k)$ is a special Cartan pair
(see Definition \ref{def:Cartan-pair}) and $B_k\subset U_j$
for some $j=j(k)$.  
To prove the theorem, it therefore suffices to show that for 
every special Cartan pair $(A,B)$ we can approximate any map 
$f\in \Lambda^r_{\Oscr} (A,Y)$ as closely 
as desired in the $\Lambda^r$ topology 
by maps $\tilde f\in \Lambda^r_{\Oscr} (D,Y)$, 
where $D=A\cup B$; the theorem then follows by a finite induction
using the sequence \eqref{eq:sequenceA}.

We first consider the case of functions, that is, $Y=\C$. 
Fix such a pair $(A,B)$ and a function 
$f\in \Lambda^r_\Oscr(A)$. Set $C=A\cap B$.
We can find a holomorphic function $g$ on a neighbourhood of $B$ 
which approximates $f|_C$ as closely as desired in $\Lambda^r(C)$. 
To do this, we proceed as follows.

By the assumption, there is a neighbourhood $W\subset X$ 
of $B$ and a holomorphic coordinate map $\psi:W\to \wt W\subset\C^n$,
with $n=\dim X$, such that the sets $\wt C=\psi(C)$ and $\wt B=\psi(B)$
are strongly convex. Choose a point $p$ in the interior of $\wt C$;
we may assume that $p=0\in\C^n$. For $t\in (0,1)$
the holomorphic map $\phi_t:W\to W$ defined by 
$\phi_t(x)=\psi^{-1}(t \psi(x))$, $x\in W$, satisfies
$\phi_t(C)\subset \mathring C$. 
The function $f_t = f \circ\phi_t$ is then holomorphic on a 
neighbourhood $V_t \subset W$ of $C$ and it approximates
$f|_C$ in $\Lambda^r(C)$ for $t$ close to $1$. 
Fix $t$ and pick a compact neighbourhood 
$C'\subset V_t$ of $C$ such that $\psi(C')$ is convex. 
By the Oka--Weil theorem, we can approximate $f_t$ uniformly on 
$C'$ by a function $g\in \Oscr(B)$. By Cauchy estimates, 
this gives approximation of $f_t|_C$ by $g|_C$ in $\Lambda^r(C)$. 
If the approximation is close enough in both steps then $g$ 
approximates $f|_C$ to the desired precision in $\Lambda^r(C)$. 

Condition (ii) in Definition \ref{def:Cartan-pair}
ensures the existence of a smooth function $\chi: X \to[0,1]$ 
which equals $1$ on a neighbourhood of $\overline {A\setminus B}$ 
and equals $0$ on a neighbourhood of $\overline {B \setminus A}$. Set
\[
	u=\chi f+(1-\chi)g\in  \Lambda^r(D).
\] 
Note that $u=f$ on $\overline{A\setminus B}$, $u=g$ on 
$\overline{B\setminus A}$, $f-u = (1-\chi) (f-g)$ on $A$, and hence 
\begin{equation}\label{eq:est1}
	\|f-u\|_{\Lambda^r(A)}= \|(1-\chi) (f-g)\|_{\Lambda^r(A)}
	\le c_0 \|f-g\|_{\Lambda^r(C)}
\end{equation}
for some $c_0>0$ depending only on $\chi$.
Furthermore, $\dibar u = (f-g)\dibar\chi \in \Lambda^r_{0,1}(D)$ 
is a $\dibar$-closed $(0,1)$-form whose support is disjoint from 
$\overline{A\setminus B}\cup \overline{B\setminus A}$. It follows that 
\[ 
	\|\dibar u\|_{\Lambda^r_{0,1}(D)} \le c_1 \|f-g\|_{\Lambda^r(C)}
\] 
for some $c_1>0$ depending on $\chi$. 
By Theorem \ref{th:dibar} there exists $\tilde u\in \Lambda^r(D)$
satisfying $\dibar \tilde u=\dibar u$ and 
\begin{equation}\label{eq:est3}
	\|\tilde u\|_{\Lambda^r(D)}\le 
	c_2 \|\dibar u\|_{\Lambda^r(D)}
	\le c_1c_2 \|f-g\|_{\Lambda^r(C)}
\end{equation}
for some $c_2>0$ depending only on $D$.
The function $\tilde f=u-\tilde u \in \Lambda^r(D)$ then 
satisfies $\dibar \tilde f=0$, so it is holomorphic on $\mathring D$. 
On $A$, we have $f-\tilde f = (f - u) + \tilde u$, and it  
follows from \eqref{eq:est1}--\eqref{eq:est3} that 
\[
	\|f - \tilde f\|_{\Lambda^r(A)}\le (c_0+c_1c_2)\|f-g\|_{\Lambda^r(C)},
\]
which can be made arbitrarily small by a suitable choice of $g$.
This proves the theorem for functions.

Assume now that $Y$ is a complex manifold.
Every map $f\in \Lambda^r_\Oscr(\overline \Omega,Y)$ for $r>0$ 
is continuous on $\overline\Omega$ and holomorphic on $\Omega$, 
so its graph has a Stein neighbourhood in $X\times Y$
(see \cite[Theorem 1.2]{Forstneric2007AJM} or 
\cite[Corollary 8.11.2]{Forstneric2017E}). The proof is then reduced to 
the case of functions as in \cite[Theorem 8.11.4]{Forstneric2017E},
using the fact that postcompositions by holomorphic maps
preserve the Zygmund spaces.
\end{proof}

%
%  PARAMETRIC VERSION OF THEOREM ON APPROXIMATION
%
%
Theorem \ref{th:approximation} has the following parametric extension.

%
%  PARAMETRIC APPROXIMATION THEOREM
%
\begin{theorem}\label{th:approximation-par}
Assume that $\Omega$ is a relatively compact, smoothly bounded, strongly 
pseudoconvex domain in a Stein manifold $X$, $Y$ is a complex manifold,
and $P$ is a compact Hausdorff space. Every continuous map 
$f:P\to \Lambda^r_\Oscr(\overline \Omega,Y)$, $r>0$, 
can be approximated in the $\Lambda^r(\overline \Omega,Y)$ topology
by continuous maps $\tilde f: P \to \Oscr(\overline \Omega,Y)$.
If in addition $Q$ is a closed subspace of $P$ which is a 
strong neighbourhood deformation retract and 
$f|_{Q}:Q \to \Oscr(\overline \Omega,Y)$, 
we can choose $\tilde f$ to agree with $f$ on $Q$.
\end{theorem}

\begin{proof}
Denote by $pr_X:X\times Y\to X$ the projection onto the first factor. 
Fix a point $p_0\in P$. 
The map $f(p_0)\in \Lambda^r_\Oscr(\overline \Omega,Y)$
is continuous on $\overline \Omega$, so its graph over $\overline \Omega$
has an open Stein neighbourhood $V_0 \subset X\times Y$. 
Furthermore, $V_0$ can be chosen to be fibrewise biholomorphic 
(with respect to the projection $pr_X$) 
to a domain with convex fibres in a holomorphic vector bundle
$E_0 \to U_0$ over an open neighbourhood $U_0\subset X$ of 
$\overline \Omega$ (see \cite[Theorem 1.2]{Forstneric2007AJM} or
\cite[Theorem 8.11.1]{Forstneric2017E}).  
With respect to such a biholomorphism, the notion of a convex combination
of points in the fibres of $pr_X:V_0\to X$ is well defined. 
By Theorem \ref{th:approximation} there is  
a holomorphic map $\tilde f(p_0) \in \Oscr(U_0,Y)$  
on a neighbourhood $U_0\subset X$ of $\overline \Omega$ which 
approximates $f(p_0)$ in $\Lambda^r(\overline \Omega,Y)$ 
to a desired precision and whose graph is contained in $V_0$. 
For $p\in P$ close to $p_0$, the graph of 
$f(p)\in \Lambda^r_\Oscr(\overline \Omega,Y)$ lies 
in $V_0$ and $\tilde f(p_0)|_{\overline\Omega}$ is an approximant to $f(p)$. 
Repeating the same procedure at other points of $P$ gives a finite open 
covering $\Pcal=\{P_i\}_{i=0}^m$ of $P$ and for each $i=0,\ldots,m$ 
a pair of open Stein domains $U_i\subset X$, $V_i\subset X\times Y$ and 
a map $\tilde f_i\in \Oscr(U_i,Y)$ such that the graph of every map
$f(p)$, $p\in U_i$, and also of $\tilde f_i$, 
lies in $V_i$, and $f(p)$ is close to $\tilde f_i$ to a desired 
precision in $\Lambda^r(\overline \Omega,Y)$ for all $p\in U_i$. 
In view of the fibrewise convex structure of every domain $V_i$,
we can use the {\em method of successive patching} 
(see \cite[p.\ 78, p.\ 282]{Forstneric2017E} for the details) 
to obtain a map $\tilde f:P\to \Oscr(\overline \Omega,Y)$ 
approximating $f$ to a desired precision. This proves 
the first part of the theorem.

For the last statement in the theorem, we precompose $f$ by a continuous map 
$\psi:P\to P$ which maps a small neighbourhood
$Q'_0\subset P$ of $Q$ to itself, it retracts a neighbourhood
$Q_0\subset Q'_0$ of $Q$ onto $Q$, and it equals the identity map on 
$P\setminus Q'_0$. (Such $\psi$ exists since $Q$ is a strong deformation 
neighbourhood retract in $P$.) This yields a continuous map 
$f_0=f\circ \psi : P\to \Lambda^r_\Oscr(\overline \Omega,Y)$
which approximates $f$, it agrees with $f$ on $Q$, and such that 
$f_0|_{Q_0}:Q_0 \to \Oscr(\overline \Omega,Y)$. Applying the 
above procedure to $f_0$ yields a map 
$\tilde f:P\to \Oscr(\overline \Omega,Y)$ approximating
$f$ which agrees with $f$ on $Q$.
\end{proof}


%
%
%	SECTION: CARTAN'S LEMMA AND THEOREM A
%
%
\section{Cartan's lemma and Theorem A for vector bundles of class 
$\Lambda^r_\Oscr$}
\label{sec:vectorbundles}

The notion of a vector bundle of Zygmund class 
$\Lambda^r_\Oscr(\overline\Omega)$, $r>0$, 
was introduced in Subsect.\ \ref{ss:Lr-bundles}. 
In this section, we prove the following version of Cartan's Theorem A for 
such vector bundles. 

%
%  EXISTENCE OF GLOBAL SECTIONS OF VECTOR BUNDLES
%
\begin{theorem}\label{th:TheoremA}
Let $\overline \Omega$ be a compact, smoothly bounded, 
strongly pseudoconvex domain in a Stein manifold $X$,  
and let $\pi:E\to \overline \Omega$ be a vector bundle 
of class $\Lambda^r_\Oscr(\overline\Omega)$ for some $r>0$. 
There exist finitely many sections in $\Gamma^r_\Oscr(\overline\Omega,E)$ 
spanning every fibre $E_x:=\pi^{-1}(x)$ for $x\in \overline \Omega$.
\end{theorem}

The classical Theorem A of Henry Cartan gives such a statement 
for holomorphic vector bundles on finite dimensional Stein spaces.
For vector bundles of class $\Ascr^r(\overline\Omega)$ with $r\in\Z_+$,
and for more general coherent analytic sheaves of this class,
the analogue of Theorem \ref{th:TheoremA} is due to 
Leiterer \cite{Leiterer1976II} for domains in Euclidean spaces and
to Heumenann \cite[Theorems 2, 6]{Heunemann1986II} for domains
in Stein manifolds.

\begin{remark}\label{rem:TheoremA}
An equivalent statement of Theorem \ref{th:TheoremA} 
is that every vector bundle $E\to \overline \Omega$ as in 
the theorem admits a vector bundle epimorphism
$\Phi:\overline\Omega\times \C^m \to E$ of class
$\Lambda^r_\Oscr(\overline\Omega)$ for some $m\in\N$.
Indeed, if $\xi_1,\ldots,\xi_m:\overline\Omega\to E$
are sections of class $\Lambda^r_\Oscr(\overline\Omega)$ 
spanning every fibre of $E$, then the map 
$\Phi:\overline\Omega\times \C^m \to E$ given by 
\begin{equation}\label{eq:Phi}
	\Phi(x,z_1,\ldots,z_m)= \sum_{i=1}^m z_i\xi_i(x) \in E_x=\pi^{-1}(x),
	\quad x\in\overline\Omega,\ z=(z_,\ldots,z_m)\in\C^m
\end{equation}
is a vector bundle epimorphism of class 
$\Lambda^r_\Oscr(\overline\Omega)$. Conversely,
given such an epimorphism, the images of standard basis sections
of the trivial bundle $\overline\Omega\times \C^m$ generate each fibre of $E$. 
\end{remark}


We also have the following embedding result for vector bundles
of class $\Lambda^r_\Oscr(\overline\Omega)$. 

\begin{corollary}\label{cor:embedding}
Given a vector bundle $\pi:E\to \overline \Omega$ of class
$\Lambda^r_\Oscr(\overline\Omega)$ as in Theorem \ref{th:TheoremA},
there exists a vector bundle embedding
$E\hookrightarrow \overline \Omega\times \C^m$ of class
$\Lambda^r_\Oscr(\overline\Omega)$ for some $m\in\N$.
\end{corollary}

\begin{proof}
Theorem \ref{th:TheoremA} gives a vector bundle epimorphism
$\Phi: \overline \Omega \times \C^m \to E^*$ of class
$\Lambda^r_\Oscr(\overline\Omega)$. Its dual $\Phi^*:(E^*)^*=E \to 
(\overline \Omega \times \C^m)^* \cong \overline \Omega \times \C^m$
is a vector bundle embedding of class 
$\Lambda^r_\Oscr(\overline\Omega)$.
\end{proof}

Denote by $M_n=M_n(\C)\cong \C^{n\times n}$ 
the space of complex $n\times n$ matrices
and note that $GL_n=GL_n(\C)\subset M_n$. 
By $I\in GL_n$ we denote the identity matrix. 
In the proof of Theorem \ref{th:TheoremA} we shall use the following
version of Cartan's lemma. See also the parametric version,
Lemma \ref{lem:Cartan-par}, and \cite{Cartan1940}
or \cite[Section VI.\ E]{GunningRossi1965} for the classical 
Cartan's lemma.

%
%  CARTAN'S SPLITTING LEMMA
%
\begin{lemma}\label{lem:Cartan}
Let $(A,B)$ be a Cartan pair in a complex manifold $X$
(see Definition \ref{def:Cartan-pair}) such that $C:=A\cap B$
is holomorphically convex in $B$.
Given a map $\gamma\in \Lambda^r_\Oscr(C,GL_n)$ 
which is homotopic to the constant map $C\ni x\mapsto I\in GL_n$ 
by a path in $\Cscr(C,GL_n)$ and a number $\epsilon>0$, there are 
maps $\alpha \in \Lambda^r_\Oscr(A,GL_n)$ and 
$\beta \in \Lambda^r_\Oscr(B,GL_n)$ such that 
$\|\alpha-I\|_{\Lambda^r(A)} <\epsilon$ and 
\begin{equation}\label{eq:split}
	\gamma = \alpha^{-1} \cdotp \beta\ \ \ \text{holds on $C$}.
\end{equation}
\end{lemma}

\begin{proof}
We first construct the product splitting \eqref{eq:split}
for $\gamma$ close to $I\in GL_n$. Write 
$\gamma=I+c$ with $c\in \Lambda^r_\Oscr(C,M_n)$.
We have the following solution to the Cousin problem 
in $\Lambda^r_\Oscr$.

%
%  Cousin-I
%
\begin{lemma}\label{lem:Cousin}
There are bounded linear operators 
$\Acal:\Lambda^r_\Oscr(C)\to \Lambda^r_\Oscr(A)$,
$\Bcal:\Lambda^r_\Oscr(C)\to \Lambda^r_\Oscr(B)$ with
\begin{equation}\label{eq:AplusB}
	\Acal c - \Bcal c =c\ \ \ \text{for all}\ c\in \Lambda^r_\Oscr(C).
\end{equation}
\end{lemma}

\begin{proof}
Set $D=A\cup B$. Condition (ii) in Definition \ref{def:Cartan-pair}
implies that there is a smooth function $\chi: X \to[0,1]$ 
which equals $0$ on a neighbourhood of $\overline {A\setminus B}$ 
and equals $1$ on a neighbourhood of $\overline {B \setminus A}$. 
Given $c\in \Lambda^r_\Oscr(C)$, we have $c\chi \in \Lambda^r(A)$,
$c(1-\chi)\in \Lambda^r(B)$, and 
$\dibar (c \chi)=c\, \dibar \chi \in \Lambda^r_{0,1}(D)$ 
is a closed $(0,1)$-form on $D$ of class $\Lambda^r(D)$ 
supported on $C$. Let 
\[
	T=\dibar^* N: \{\omega\in \Lambda^r_{0,1}(D):\dibar \omega=0\}
	\to \Lambda^r(D)
\]
denote the (bounded, linear) canonical solution operator 
for the $\dibar$-equation for $(0,1)$-forms of class
$\Lambda^r(D)$; see Theorem \ref{th:dibar}. 
The linear operators on $\Lambda^r_\Oscr(C)$ defined by 
\[
	\Acal c= c \, \chi - T(c\, \dibar \chi)\in \Lambda^r_\Oscr(A), \quad 
	\Bcal c= c(\chi-1) - T(c\, \dibar \chi) \in \Lambda^r_\Oscr(B)
\] 
for all $c\in \Lambda^r_\Oscr(C)$ then satisfy the lemma.
\end{proof}

Applying Lemma \ref{lem:Cousin} componentwise gives
bounded linear operators 
\[
	\Acal:\Lambda^r_\Oscr(C,M_n) \to \Lambda^r_\Oscr(A,M_n),
	\quad 
	\Bcal:\Lambda^r_\Oscr(C,M_n) \to \Lambda^r_\Oscr(B,M_n)
\]
satisfying $\Acal-\Bcal=\Id$ on $\Lambda^r_\Oscr(C,M_n)$;
see \eqref{eq:AplusB}. We define a map 
$\Phi:U\to \Lambda^r_\Oscr(C,GL_n)$ on a small neighbourhood 
$U\subset \Lambda^r_\Oscr(C,M_n)$ of $c=0$ by
\[
	\Phi(c) = (I-\Acal c)^{-1} (I-\Bcal c)
	\in \Lambda^r_\Oscr(C,GL_n), \quad c\in U.
\]
Clearly, $\Phi$ is smooth and satisfies $\Phi(0)=I$ 
and $d\Phi_0 = \Acal - \Bcal =\Id$. It follows that 
$\Phi$ has a smooth right inverse $\Psi$ on a neighbourhood 
$V\subset \Lambda^r_\Oscr(C,GL_n)$ of $I$ such that $\Psi(I)=0$ and 
\[
	(\Phi\circ \Psi)(\gamma) = 
	\left(I-\Acal \circ \Psi(\gamma)\right)^{-1}  
	\left(I-\Bcal \circ\Psi(\gamma)\right) = \gamma 
	\quad \text{for all}\ \gamma\in V.
\]
The smooth operators $\wt \Acal:V\to \Lambda^r_\Oscr(A,GL_n)$,
$\wt \Bcal:V\to \Lambda^r_\Oscr(B,GL_n)$ defined by
\[
	\wt \Acal = I - \Acal \circ\Psi,\qquad 
	\wt \Bcal = I - \Bcal \circ\Psi
\] 
then provide a splitting 
$\gamma=(\wt \Acal \gamma)^{-1} (\wt \Bcal \gamma)$ 
for $\gamma\in V$ (see \eqref{eq:split}) satisfying 
\begin{equation}\label{eq:estAB}
	\|\wt \Acal \gamma-I\|_{\Lambda^r(A)} \le const 
	\|\gamma-I\|_{\Lambda^r(C)},
	\quad 
	\|\wt \Bcal \gamma-I\|_{\Lambda^r(B)} \le const 
	\|\gamma-I\|_{\Lambda^r(C)}
\end{equation}
for some $const>0$ depending only on $(A,B)$ and $r$.

This proves the lemma for maps $\gamma\in \Lambda^r_\Oscr(C,GL_n)$
near the constant map $C\ni x\mapsto I$.
The general case is obtained as follows.
By Theorem \ref{th:approximation} we can approximate 
$\gamma$ as closely as desired in $\Lambda^r_\Oscr(C)$ by 
a holomorphic map ${\gamma\,}' :U\to GL_n$ from an 
open neighbourhood $U\subset X$ of $C$. Since
$GL_n$ is an Oka manifold (every complex homogeneous
manifold is an Oka manifold by Grauert 
\cite{Grauert1957I,Grauert1957II}; see also 
\cite[Proposition 5.6.1]{Forstneric2017E}), 
$\gamma$ is homotopic to the constant map,
and $C$ is holomorphically convex in $B$, 
the Oka principle \cite[Corollary 5.4.5]{Forstneric2017E}
shows that ${\gamma\,}'$ can be approximated uniformly
on a compact neighbourhood $C'\subset U$ of $C$ by a holomorphic
map $\tilde \gamma\in \Oscr(B,GL_n)$. 
Then, $\gamma= (\gamma {\tilde \gamma}^{-1}) \tilde\gamma$
on $C$, and $\gamma {\tilde \gamma}^{-1}$ is close to $I$ 
in $\Lambda^r(C,GL_n)$. By the first part, we have
$\gamma {\tilde \gamma}^{-1}=\alpha^{-1}\tilde \beta$
with $\alpha \in \Lambda^r_\Oscr(A,GL_n)$ close to $I$ and 
$\tilde \beta \in \Lambda^r_\Oscr(B,GL_n)$.
Setting $\beta=\tilde\beta \tilde\gamma\in \Lambda^r_\Oscr(B,GL_n)$ 
gives $\gamma=\alpha^{-1}\beta$ as in \eqref{eq:split}.
\end{proof}

%
% 	CARTAN LEMMA WITH PARAMETERS
%
Lemma \ref{lem:Cartan} has the following generalisation
to the parametric case.

\begin{lemma}\label{lem:Cartan-par}
Assume that $X$ and $(A,B)$ are as in Lemma  \ref{lem:Cartan}.
Given a compact Hausdorff space $P$, a closed subspace 
$Q\subset P$ which is a strong neighbourhood deformation retract, 
and a continuous map 
$\gamma : P\to \Lambda^r_\Oscr(C,GL_n)$ such that $\gamma(p)=I$
for all $p\in Q$ and there is a homotopy 
$\gamma_t:P\to \Cscr(C,GL_n)$ $(t\in [0,1])$ which is fixed 
on $Q$ such that $\gamma_0=I$ and $\gamma_1=\gamma$,
there are continuous maps 
\[
	\alpha:P \to \Lambda^r_\Oscr(A,GL_n),\qquad 
	\beta:P\to  \Lambda^r_\Oscr(B,GL_n)
\]
such that $\|\alpha-I\|_{\Lambda^r(A)}$ is arbitrarily small,
$\alpha|_Q=\beta|_Q=I$, and 
$\gamma = \alpha^{-1} \cdotp \beta$ holds on $C$.
\end{lemma}

\begin{proof}
By Theorem \ref{th:approximation-par} and the argument
in the last part of the proof of Lemma \ref{lem:Cartan}, we can 
reduce to the case when $\gamma$ is close to the 
constant map $P\to I\in GL_n$. Since the splitting 
\eqref{eq:AplusB} in Lemma \ref{lem:Cousin} is given by bounded 
linear operators, it also applies to the parametric case. The proof
of the first part of Lemma \ref{lem:Cartan} then carries over verbatim.
\end{proof}


\begin{remark} \label{rem:Liegroup} 
Lemmas \ref{lem:Cartan} and \ref{lem:Cartan-par} also hold, with the
same proofs, for maps with values in any complex Lie group $G$ in place
of $GL_n(\C)$. In this case, we replace the matrix algebra
$M_n$ (which is the Lie algebra of $GL_n(\C)$)
by the Lie algebra of $G$.
Furthermore, the analogous results hold with the spaces
$\Lambda^r_\Oscr(\overline\Omega,G)$ replaced by any space
$\Fcal_\Oscr(\overline\Omega,G)$ described in Remark \ref{rem:Fcal}.
\end{remark}

The main step in the proof of Theorem \ref{th:TheoremA} is 
given by the following lemma.

%
% EXTENSION OF A VECTOR BUNDLE EPIMORPHISM TO A BUMP
%
\begin{lemma}\label{lem:ext-to-bump}
Assume that $(A,B)$ is a special Cartan pair 
(see Definition \ref{def:Cartan-pair}).
Let $\pi:E\to D=A\cup B$ be a vector bundle of class
$\Lambda_\Oscr^r(D)$, $r>0$, such that $E|_B$ is 
$\Lambda_\Oscr^r(B)$-isomorphic to a trivial bundle.
Then, every vector bundle epimorphism 
$\Phi:A\times\C^n \to E|_A$ of class $\Lambda_\Oscr^r(A)$
can be approximated in $\Lambda^r(A)$ by vector bundle 
epimorphisms $\wt \Phi:D\times\C^n \to E$ of class $\Lambda_\Oscr^r(D)$. 
\end{lemma}

\begin{proof}
Let $e_1,\ldots,e_n$ be sections of $E|_A$ of class $\Lambda^r_\Oscr(A)$
which are $\Phi$-images of the standard basis sections
of $A\times\C^n$. Also, let $g_1,\ldots,g_m$ with $m=\rank E$
be basis sections of class $\Lambda^r_\Oscr(B)$ 
of the trivial bundle $E|_B\cong B\times\C^m$. On $C$, we have 
$e_i=\sum_{j=1}^m \gamma_{i,j}g_j$ ($i=1,\ldots,n$) 
with $\gamma_{i,j}\in\Lambda^r_\Oscr(C)$, and 
the $n\times m$ matrix $\Gamma'(x)=(\gamma_{i,j}(x))$ 
has rank $m$ at every point $x\in C$. We claim that there is
an $n\times n$ matrix function 
\[
	\Gamma=(\gamma_{i,j})_{i,j=1}^n =(\Gamma',\Gamma'') 
	\in \Lambda^r_\Oscr(C,GL_n)
\]
whose left hand side $n\times m$ submatrix equals $\Gamma'$. 
This is equivalent to saying that the subbundle $E'$ of  
$C\times\C^n$, spanned by the columns of $\Gamma'$, 
has a trivial complementary subbundle $E''\subset C\times\C^n$ of class 
$\Lambda^r_\Oscr(C)$. The existence of a complementary subbundle
$\wt E''$ of class $\Ascr(C)$ (continuous on $C$ and holomorphic 
on $\mathring C$) follows from \cite[Theorem 3]{Heunemann1986I}.
Since $C$ is contractible, $\wt E''$ is trivial as a bundle of class 
$\Ascr(C)$ by \cite[Theorem 2]{Heunemann1986I}. 
Let $\wt \Gamma''$ be an $n\times (n-m)$ matrix of class $\Ascr(C)$
spanned by the basis sections of $\wt E''$. 
Since $C$ is convex, we can approximate $\wt \Gamma''$ uniformly on $C$ 
by a matrix $\Gamma''$ of class $\Lambda^r_\Oscr(C)$. If the approximation 
is close enough then 
$\Gamma=(\Gamma',\Gamma'') \in \Lambda^r_\Oscr(C,GL_n)$. 
Set $g_{m+1}=0,\ldots,g_n=0$. It follows that 
\[
	(e_1,e_2,\ldots,e_n)^t = \Gamma (g_1,g_2,\ldots,g_n)^t,  
\]
where the superscript $t$ denotes the transpose. By 
Lemma \ref{lem:Cartan} we have $\Gamma=\alpha^{-1}\beta$
where $\alpha\in \Lambda^r_\Oscr(A,GL_n)$ is close to $I$
and $\beta \in \Lambda^r_\Oscr(B,GL_n)$. Then,
\[
	\alpha (e_1,e_2,\ldots,e_n)^t = \beta (g_1,g_2,\ldots,g_n)^t
\]
holds on $C$. The resulting collection on $n$ sections of $E\to D$ 
of class $\Lambda^r_\Oscr(D)$ determines a vector bundle epimorphism 
$\wt \Phi:D\times \C^n\to E$ of class $\Lambda^r_\Oscr(D)$ 
(see \eqref{eq:Phi}) which approximates $\Phi$ in $\Lambda^r_\Oscr(A)$.
\end{proof}


\begin{proof}[Proof of Theorem \ref{th:TheoremA}]
Pick a smooth strongly plurisubharmonic function $\rho$ on a 
neighbourhood $U$ of $\overline\Omega$
such that $\Omega=\{x\in U:\rho(x)<0\}$ and $d\rho_x\ne 0$
for every point $x\in b\Omega =\{\rho=0\}$. Choose $c<0$ such that
$\rho$ has no critical values on $[c,0]$ and 
set $A_0=\{\rho\le c\}$. By \cite[Lemma 5.10.3]{Forstneric2017E}, 
given an open cover $\Ucal=\{U_j\}$ of 
$\overline{\Omega\setminus A_0}=\{c\le \rho\le 0\}$ consisting of
holomorphic coordinate charts $U_j\subset X$,   
there are compact, smoothly bounded, strongly pseudoconvex domains 
$A_0\subset A_1 \subset \cdots \subset A_{k_0}=\overline\Omega$
for some $k_0\in\N$ such that for every $k=0,1,\ldots,k_0-1$ we have 
$A_{k+1}=A_k\cup B_k$, where $(A_k,B_k)$ is a special Cartan pair
(see Definition \ref{def:Cartan-pair}) and $B_k\subset U_j$
for some $j=j(k)$. We choose the sets $U_j$ small enough so that 
the restricted bundle $E|_{B_k}$ is trivial for $k=0,1,\ldots,k_0-1$.
Since $E$ is holomorphic over $\Omega$,
there is a holomorphic vector bundle epimorphism
$\Phi_0:A_0\times \C^n \to E|_{A_0}$ for some $n\in \N$. 
We inductively apply Lemma \ref{lem:ext-to-bump}
to find vector bundle epimorphisms
$\Phi_k:A_k\times \C^n \to E|_{A_k}$ of class $\Lambda^r_\Oscr(A_k)$
for $k=1,\ldots,k_0$ such that $\Phi_{k+1}$ approximates $\Phi_k$ 
on $A_k$ for every $k=0,1,\ldots,k_0-1$.
Then, $\Phi:=\Phi_{k_0}:\overline \Omega \times \C^n\to E$ 
satisfies the conclusion of the theorem.
\end{proof}

Next, we show that every complex vector subbundle 
of class $\Lambda^r_\Oscr(\overline\Omega)$ is complemented.

%
%  COMPLEMENTARY SUBBUNDLE
%
\begin{theorem}\label{th:complement}
Let $\pi:E\to \overline\Omega$ be a complex vector bundle of class
$\Lambda^r_\Oscr(\overline\Omega)$. For every 
complex vector subbundle $E'\subset E$ of class 
$\Lambda^r_\Oscr(\overline\Omega)$ there is a complementary 
to $E'$ complex vector subbundle $E''\subset E$ of class 
$\Lambda^r_\Oscr(\overline\Omega)$ such that $E\cong E'\oplus E''$
as complex vector subbundles of class $\Lambda^r_\Oscr(\overline\Omega)$.
\end{theorem}

\begin{proof}
We first consider the case when $E= \overline\Omega\times\C^n$
is a trivial bundle. By \cite[Theorem 3]{Heunemann1986II}
there is a vector subbundle $\wt E\subset E$ of class $\Ascr(\overline\Omega)$ 
complementary to $E'$. To complete the proof,
it suffices to approximate $\wt E$ sufficiently closely by a subbundle 
$E''\subset E=\overline\Omega\times\C^n$ of class 
$\Lambda^r_\Oscr(\overline\Omega)$. We follow the idea
in \cite[Proof of Lemma 2.2]{ForstnericLowOvrelid2001}.

Let $\wt L: \overline \Omega \to M_n=\Lin_\C(\C^n,\C^n)$ be the unique map
such that $\wt L(x):\C^n\to\C^n$ is the projection onto $\wt E_x$ 
with kernel $E'_x$ for every $x\in \overline \Omega$.
Clearly, $\wt L$ is of class $\Ascr(\overline\Omega)$, and 
the eigenvalues of $\wt L(x)$ are $0$ and $1$ for every 
$x\in \overline\Omega$. Let $C=\{z\in \C : |z-1|=1/2\}$. 
We can approximate $\wt L$ uniformly on $\overline\Omega$ 
by a holomorphic map $L:U \to M_n$ 
on a neighbourhood $U$ of $\overline\Omega$. 
Assuming that the approximation is close enough, 
$L(x)$ has no eigenvalues on the curve $C$ for $x\in\overline\Omega$,
we have $\C^n=V_{x,+} \oplus V_{x,-}$
where $V_{x,+}$ resp.\ $V_{x,-}$ are the $L(x)$-invariant 
subspaces of $\C^n$ spanned by the generalised
eigenvectors of $L(x)$ inside resp.\ outside of $C$, and  
$V_{x,+}$ is close to the subspace $\wt E_x=\wt L(x)(\C^n) \subset\C^n$. 
The map 
\[
	P(L(x))= \frac{1}{2\pi \imath} 
	\int_{\zeta\in C} \left(\zeta I-L(x) \right)^{-1}\, d\zeta \in M_n
\]
is the projection of $\C^n$ onto $V_{x,+}$ with kernel $V_{x,-}$ 
(see \cite{GohbergLancasterRodman1986}), it 
depends holomorphically on $x$ in a neighbourhood of $\overline\Omega$,
and it approximates $L(x)$ uniformly on $x\in \overline\Omega$. 
Assuming that the approximations are close enough, the subbundle 
$E''\subset E$ with fibres $E''_x=V_{x,+}=P(L(x))(\C^n)$, 
$x\in \overline\Omega$, is complementary to $E'$, 
and $E''$ is holomorphic on a neighbourhood of $\overline\Omega$. 

For a general vector bundle $E\to\overline\Omega$
of class $\Lambda^r_\Oscr(\overline\Omega)$, we apply 
Theorem \ref{th:TheoremA} to find a vector bundle
epimorphism $\Phi:\overline\Omega\times\C^N\to E$ of class
$\Lambda^r_\Oscr(\overline\Omega)$. The preimage $\Phi^{-1}(E')$
is a vector subbundle of $\overline\Omega\times\C^N$
class $\Lambda^r_\Oscr(\overline\Omega)$, 
so it has a complementary subbundle 
$\wt E''\subset \overline\Omega\times\C^N$
of class $\Lambda^r_\Oscr(\overline\Omega)$. 
Its image $E''=\Phi(\wt E'')$ is then a subbundle of $E$
of class $\Lambda^r_\Oscr(\overline\Omega)$ which is 
complementary to $E'$.
\end{proof}



%
%  LEITERER'S STABILITY THEOREM
%
\begin{remark}\label{rem:stability}
In the sequel, we shall be using the notion of a continuous family 
of holomorphic vector bundles on a smoothly bounded domain 
$\Omega\Subset X$ which are continuous or better on $\overline\Omega$;
see Leiterer \cite[Definition 2.14, p.\ 70]{Leiterer1990} (where
the parameter space is $[0,1]$) or the statement of his stability 
theorem \cite[Theorem 2.7]{Leiterer1990} for the general case.
This means that, locally in the parameter, the family is defined
by a continuous family of 1-cocycles on the same finite open covering
of $\overline\Omega$. The stability theorem says that, 
under suitable cohomological conditions on the endomorphism bundle 
$\mathrm{Ad}(E)$ of a vector bundle $E\to \overline\Omega$ of class 
$\Ascr(\overline\Omega)$ (see the statement of 
\cite[Theorem 2.7]{Leiterer1990}), all nearby vector bundles of the same
class are isomorphic to $E$ with a continuous dependence on the
parameter. This follows from an implicit function theorem
in Banach spaces given by \cite[Theorem 2.9]{Leiterer1990}.
The aforementioned cohomological conditions 
are satisfied if for every continuous $\dibar$-closed 
$E$-valued $(0,q)$-form $\phi$ on $\overline\Omega$,
$q>0$, the equation $\dibar \psi=\phi$ can be solved with a continuous
$\psi$ on $\overline\Omega$ (see \cite[proof of Theorem 2.12]{Leiterer1990}). 
This holds in particular if $\Omega$ is strongly pseudoconvex; 
see \cite[Lemma 2.13]{Leiterer1990}. By Theorems \ref{th:dibar}
and \ref{th:complement}, the same conclusion holds in 
H\"older--Zygmund spaces, so the analogue of 
\cite[Theorem 2.7]{Leiterer1990} also holds for vector bundles
of class $\Lambda^r_\Oscr(\overline\Omega)$. 
For later reference we state this result explicitly.

\begin{theorem}\label{th:stability}
Assume that $\overline \Omega$ is a compact, smoothly bounded, 
strongly pseudoconvex domain in a Stein manifold $X$,  
$P$ is a topological space, and $E_p\to \overline \Omega$ $(p\in P)$ 
is a continuous family of vector bundles of class 
$\Lambda^r_\Oscr(\overline\Omega)$, $r>0$. 
Then for each $p_0\in P$ there are a neighbourhood $P_0\subset P$
and a continuous family of vector bundle isomorphisms
$\phi_p:E_p \stackrel{\cong}{\longrightarrow} E_{p_0}$ $(p\in P_0)$ 
of class $\Lambda^r_\Oscr(\overline\Omega)$. 
\end{theorem}
\end{remark}



%
%  PARAMETRIC THEOREM A
%
By using this result, we obtain the following parametric version of Theorem \ref{th:TheoremA}. 

\begin{theorem}
\label{th:TheoremApar}
Assume that $\overline \Omega$ is a compact, smoothly bounded, 
strongly pseudoconvex domain in a Stein manifold $X$,  
$P$ is a compact Hausdorff space, 
and $E_p\to \overline \Omega$ $(p\in P)$ is a continuous
family of vector bundles of class 
$\Lambda^r_\Oscr(\overline\Omega)$, $r>0$. 
There exist an integer $N\in \N$ and a 
continuous family of vector bundle epimorphisms
$\Phi_p: \overline \Omega\times \C^N \to E_p$, $p\in P$,
of class $\Lambda^r_\Oscr(\overline \Omega)$.
\end{theorem}

\begin{proof}
By Theorem \ref{th:stability}, every point 
$p_0\in P$ has a neighbourhood $P_0\subset P$ and a continuous 
family of vector bundle isomorphisms $\phi_p:E_p\to E_{p_0}$ 
$(p\in P_0)$ of class $\Lambda^r_\Oscr(\overline\Omega)$.
A vector bundle epimorphism
$\Psi_{p_0}: \overline \Omega\times \C^{n_0} \to E_{p_0}$,
given by Theorem \ref{th:TheoremA}, extends to a continuous
family of such epimorphisms $\Psi_p=\phi_p^{-1}\circ \Psi_{p_0}:
\overline \Omega\times \C^{n_0} \to E_{p}$ for $p\in P_0$.
Since $P$ is compact, we obtain finitely many pairs of open 
subsets $P'_i \Subset P_i\subset P$ and 
continuous families of vector bundle 
epimorphism $\Psi^i_p:\overline \Omega\times \C^{n_i}\to E_p$ 
($p\in P_i$,  $i=1,\ldots,m$) such that $\bigcup_{i=1}^m P'_i=P$ .
For every $i=1,\ldots,m$ let $\chi_i:P\to [0,1]$ be a continuous
function such that $\chi_i=1$ on $P'_i$ and $\supp \chi_i\subset P_i$.
Let $\Phi^i_p: \overline \Omega\times \C^{n_i}\to E_p$ be 
defined by $\Phi^i_p(x,z)= \Psi^i_p(x,\chi_i(p)z)$. 
Take $N=\sum_{i=1}^m n_i$. The family 
\[
	\Phi_p=\bigoplus_{i=1}^m \Phi^i_p: 
	\overline \Omega\times \C^{N}\to E_p,
	\quad p\in P,
\]
clearly satisfies the theorem.
\end{proof}


%
%  FAMILIES OF COMPLEMENTARY VECTOR SUBBUNDLES
%
We also have the following parametric version of 
Theorem \ref{th:complement}.

\begin{theorem}\label{th:complement-par}
Assume that $\Omega$ is as in Theorem \ref{th:complement}, 
$P$ is a compact Hausdorff space, and 
$\pi_p:E_p\to \overline\Omega$ is a continuous family 
of vector bundles of class $\Lambda^r_\Oscr(\overline\Omega)$.
(See Remark \ref{rem:stability}.) 
Given a continuous family $E'_p\subset E_p$ $(p\in P)$ 
of  vector subbundles of class $\Lambda^r_\Oscr(\overline\Omega)$, 
there is a continuous family $E''_p\subset E_p$ $(p\in P)$ 
of complementary vector subbundles of class 
$\Lambda^r_\Oscr(\overline\Omega)$ such that 
\begin{equation}\label{eq:directsum} 
	E_p\cong E'_p\oplus E''_p,\quad p\in P.
\end{equation}
If in addition $Q$ is a closed subspace of $P$ which is a 
strong neighbourhood deformation retract and 
$E''_p\subset E_p$ $(p\in Q)$ is a continuous family 
of complementary to $E'_p$  vector subbundles of class 
$\Lambda^r_\Oscr(\overline\Omega)$, then the family $E''_p$
can be extended to all values $p\in P$ so that 
\eqref{eq:directsum} holds.
\end{theorem}

\begin{proof}
Assume first that the family of vector bundles 
$E_p=E\to \overline \Omega$ is independent of $p\in P$.
Consider the short exact sequence of vector bundle homomorphisms 
\begin{equation}\label{eq:SES}	
	0\lra E'_p \lra E \stackrel{\phi_p}{\lra} E/E'_p \lra 0,
\end{equation}
where $\phi_p$ is the natural quotient projection for every $p\in P$. 
Theorem \ref{th:complement} gives for every $p\in P$ a 
subbundle $E''_p\subset E$ of class $\Lambda^r_\Oscr(\overline\Omega)$
complementary to $E'_p$. Such a subbundle is the image of a 
unique vector bundle monomorphism 
$\sigma_p: E/E'_p\to E$ of the same class 
satisfying $\phi_p\circ \sigma_p=\Id$ on $E/E'_p$, and vice versa.
Such $\sigma_p$ is called a splitting of the sequence \eqref{eq:SES}.
Since $\overline \Omega$ is compact and the subbundles
$E'_p$ of $E$ vary continuously with $p$, a complementary 
to $E'_{p_0}$ subbundle $E''_{p_0}\subset E$ 
is also complementary to $E'_p$ for all $p\in P$  
near $p_0$, so it gives a continuous family of
splittings $\sigma_p:E/E'_p\to E$. Furthermore, a convex linear 
combination of splittings of $\phi_p$ in \eqref{eq:SES} 
is again a splitting. Hence, the proof follows by using a continuous
partition of unity on the parameter space $P$.

The case of a continuously varying family $E_p\to \overline \Omega$, 
$p\in P$, can be reduced to the special case as follows. 
By Theorem \ref{th:stability}, every point $p_0\in P$ has a neighbourhood
$P_0\subset P$ and a continuous family of vector bundle isomorphisms 
$\phi_p:E_p\to E_{p_0}$ $(p\in P_0)$ of class 
$\Lambda^r_\Oscr(\overline\Omega)$.
Hence, the family $E_p$ is locally constant, and hence 
our proof applies locally on $P$. It remains to patch 
the locally defined continuous families of 
splittings by a continuous partition of unity on $P$. 
\end{proof}


%
%
%	A GLUING LEMMA
%
%
\section{A gluing lemma for sprays of class 
$\Lambda^r_\Oscr$}\label{sec:gluing}

Let $X$ be a complex manifold. Given a compact 
smoothly bounded subset $K$ of $X$
and an open convex domain $0\in W\subset \C^N$ (the 
precise choice of $W$ is not important; for convenience
we may take a ball or a polydisc), we shall consider maps
$\gamma : K\times W\to K\times \C^N$ of the form
%
\begin{equation}\label{eq:gamma}
     \gamma(x,w) = \bigl(x,\psi(x,w)\bigl),\quad x\in K, \ w\in W
\end{equation}
and of class $\Lambda^r_\Oscr(K\times W)$. The domain $W$ 
will have the role of a parameter space and will be allowed  
to shrink during the construction. 
Let $\Id(x,w)=(x,w)$ denote the identity map on $X\times \C^N$. 

%
%  SPLITTING LEMMA FOR CARTAN PAIRS OF CLASS C2
%
%
The following splitting lemma is an analogue of 
\cite[Proposition 5.8.1, p.\ 235]{Forstneric2017E}, 
which pertains to function spaces $\Ascr^r(K)$
with $r\in \Z_+$. It was first proved in 
\cite[Theorem 3.2]{DrinovecForstneric2008FM} 
and \cite[Lemma 3.2]{Forstneric2007AJM}.


\begin{lemma}\label{lem:splitting}
Let $(A,B)$ be a Cartan pair in a Stein manifold $X$ 
(see Def.\ \ref{def:Cartan-pair}) and set $C=A\cap B$, $D=A\cup B$. 
Given a convex domain 
$0\in W\subset\C^N$ and a number $\epsilon \in (0,1)$, 
there is a number $\delta>0$ satisfying the following.
For every map $\gamma : C\times W \to C\times \C^N$
of the form \eqref{eq:gamma} and of class 
$\Lambda^r_\Oscr(C\times W)$ satisfying
$\|\gamma-\Id\|_{\Lambda^r(C\times W)}<\delta$ 
there exist maps
\begin{equation}\label{eq:alphabeta}
 	 \alpha_\gamma : A \times \epsilon W \to A\times \C^N, \qquad
 	 \beta_\gamma :  B \times \epsilon W \to B \times \C^N
\end{equation}
of the form \eqref{eq:gamma} and of class
$\Lambda^r_\Oscr(A\times \epsilon W)$ and 
$\Lambda^r_\Oscr(B\times \epsilon W)$, respectively, 
depending smoothly on $\gamma$,
such that $\alpha_{\Id}=\Id$, $\beta_{\Id}=\Id$, and
\begin{equation}\label{eq:splitspray}
   \gamma\circ\alpha_\gamma = \beta_\gamma
   \quad \text{holds on $C\times \epsilon W$.}
\end{equation}
If $\gamma$ agrees with $\Id$ to order $k\in\N$ along $w=0$
then so do $\alpha_\gamma$ and~$\beta_\gamma$.
The analogous result holds with a continuous dependence
of maps on a parameter $p\in P$ in a compact Hausdorff space.
\end{lemma}

\begin{proof}
We follow \cite[proof of Proposition 5.8.1]{Forstneric2017E}, 
taking into account \cite[Remark 5.8.3 (B), p.\ 238]{Forstneric2017E}
and Theorem \ref{th:dibar} in this paper. In particular, 
the use of \cite[Lemma 5.8.2, p.\ 236]{Forstneric2017E}
is replaced by Lemma \ref{lem:Cousin}. We recall the proof 
and adjust it to the H\"older--Zygmund classes.

Given a number $\epsilon\in (0,1]$, consider the Banach space
$C_{\epsilon}$ consisting of all continuous maps
\[
	\psi=(\psi_1,\ldots,\psi_N):C\times \epsilon W \to \C^N
\] 
which are holomorphic in $\mathring C\times \epsilon W$ and satisfy
\[
	\psi(\cdotp,w)\in \Lambda^r(C,\C^N)\ \ 
	\text{for all $w\in \epsilon W$},
	\qquad
	\|\psi\|_{\epsilon} = \sup_{w\in \epsilon W} 
	\sum_{i=1}^N \|\psi_i(\cdotp,w)\|_{\Lambda^r (C)} < \infty.
\]
Similarly, $A_{\epsilon}$ and $B_{\epsilon}$ denote Banach spaces
of the same kind associated to the sets $A$ and $B$, respectively. 
We may identify $\psi\in C_\epsilon$ with the bounded 
holomorphic map
\[
	\epsilon W\ni w \mapsto 
	\psi(\cdotp,w)\in  \Lambda^r_\Oscr (C,\C^N).
\]
Note that $\Lambda^r_\Oscr(C\times \epsilon W) \subset C_\epsilon$
and the inclusion is continuous. Furthermore, if $0<\epsilon'<\epsilon$ 
then the restriction map 
\begin{equation}\label{eq:restriction}
	C_\epsilon\ni \psi\mapsto \psi|_{C\times \epsilon'W} 
	\in \Lambda^r_\Oscr(C\times \epsilon' W) 
\end{equation}
is continuous; see e.g. \cite[Theorem, p.\ 63]{BesovIlinNikolskii1979}.
Let 
\[
	\Acal:\Lambda^r_\Oscr(C)\to \Lambda^r_\Oscr(A),
	\qquad  
	\Bcal:\Lambda^r_\Oscr(C)\to \Lambda^r_\Oscr(B)
\]
be the bounded linear operators in Lemma \ref{lem:Cousin}, satisfying
$\Acal c - \Bcal c=c$ for all $c\in \Lambda^r_\Oscr(C)$ 
(see \eqref{eq:AplusB}). We extend them to maps $c\in C_\epsilon$ by 
$(\Acal c)(\cdotp,w) = \Acal (c(\cdotp,w))$ for each $w\in\epsilon W$,
and likewise for $\Bcal$. This gives bounded linear operators
\[
	\Acal:C_\epsilon \to A_\epsilon,\qquad \Bcal:C_\epsilon \to B_\epsilon,
\]
such that
\[
	\Acal c - \Bcal c=c \quad\text{for every $c\in C_\epsilon$.}
\]
We are given a map 
\begin{equation}\label{eq:gamma1}
	\gamma(x,w)=(x,\psi(x,w)),\quad x\in C,\ w\in W,
\end{equation}
of the form \eqref{eq:gamma}, with $\psi\in C_{1}$ close to the map
\[
	\psi_0(x,w)=w, 
\]
and a number $0<\epsilon <1$. Given $c\in C_{\epsilon}$ near $0$, we define 
\begin{equation}\label{eq:Phi}
        \Phi(\psi,c)(x,w) =
        \psi\bigl(x,w+(\Acal c)(x,w)\bigr) - \bigl(w + (\Bcal c)(x,w)\bigr),
        \quad x\in C,\ w\in \epsilon W.
\end{equation}
Since $\Acal:C_\epsilon \to A_\epsilon$ is a bounded linear operator,
$\Phi(\psi,c)$ is well defined for any $c\in C_{\epsilon}$ 
in a neighbourhood of $0$. 
We claim that $(\psi,c)\mapsto \Phi(\psi,c)$ is a smooth map from 
an open neighborhood of $(\psi_0,0)$ in the Banach space 
$C_{1} \times C_{\epsilon}$ to $C_{\epsilon}$. 
To see this, note that $\Phi$ is affine linear in $\psi$, 
and it is continuous in $c\in C_{\epsilon}$ as a map 
to $C_{\epsilon}$. Indeed, the function $\psi$ is of the form
\[
	\psi(x,w)= \sum_{j \in \Z_+^N} \psi_j(x) w^j,\quad 
	x\in C,\ w\in W\subset\C^N, 
\]
with $\psi_j\in \Lambda^r_\Oscr(C)$ for $j\in \Z_+^N$.
Choose a number $\epsilon_1\in (\epsilon,1)$.
For $c\in C_{\epsilon}$ sufficiently near $0$, the map
\[
	C\times \epsilon W \ni (x,w) \mapsto w+(\Acal c)(x,w) \in \C^N 
\] 
belongs to $C_\epsilon$ and has range compactly 
contained in $C\times \epsilon_1 W$. For such $c$, the composition
\[
	(x,w) \mapsto \psi(x,w+(\Acal c)(x,w))
\]
lies in $C_{\epsilon}$ (since it is a convergent power series in $w$ 
with coefficients in $\Lambda^r(C)$), and it depends continuously on $c$ 
as an element of $C_{\epsilon}$. Furthermore, the partial differential 
of $\Phi(\psi,c)$ with respect to the second variable $c$ 
at $c=c_0\in C_{\epsilon}$ equals
\[
     \di_c\, \Phi(\psi,c_0) c(x,w) =
     \di_w\, \psi\bigl(x, w+(\Acal c_0)(x,w)\bigr) 
     \cdotp (\Acal c)(x,w) - (\Bcal c)(x,w).
\]
This map is again affine linear in $\psi$ and continuous in all variables.
A similar argument applies to higher order differentials of $\Phi$.
Note that 
\[
    \Phi(\psi_0,c)= \Acal(c)-\Bcal(c)=c,\quad c\in C_\epsilon,
\]
so $\di_c\, \Phi(\psi_0,0)$ is the identity map on $C_\epsilon$.
The implicit function theorem gives a smooth map 
\[
	\psi \longmapsto \Ccal(\psi) \in C_{\epsilon}, 
\]
defined in a neighborhood of $\psi_0$ in $C_1$, such that 
\[
	\Phi(\psi,\Ccal(\psi))=0 \quad {\rm and} \quad 
	\Ccal(\psi_0)=0.
\]
The maps $a_\psi$ and $b_\psi$ defined by
\begin{equation}\label{eq:apsi}
        a_\psi(x,w) = w + \Acal \circ \Ccal(\psi)(x,w), \qquad
        b_\psi(x,w) = w + \Bcal\circ \Ccal(\psi)(x,w)
\end{equation}
then satisfy $a_\psi\in A_{\epsilon}$,  $a_{\psi_0}=\psi_0$, 
$b_\psi\in B_\epsilon$, $b_{\psi_0}=\psi_0$, and
\[
        \psi\bigl(x,a_\psi(x,w)\bigr) = b_\psi(x,w),
        \quad (x,w)\in C \times \epsilon W.
\]
Furthermore, after decreasing $\epsilon>0$ slightly, the comment 
after \eqref{eq:restriction} shows that 
\[
	a_\psi\in \Lambda^r_\Oscr(A\times \epsilon W),
	\qquad 
	b_\psi\in \Lambda^r_\Oscr(B\times \epsilon W).
\]
The associated maps
\[
	\alpha_\gamma (x,w) = \bigl(x,a_\psi(x,w)\bigr),  \quad
	\beta_\gamma (x,w)  = \bigl(x,b_\psi(x,w)\bigr)
\]
depend smoothly on the map $\gamma$ \eqref{eq:gamma1} and satisfy 
condition \eqref{eq:splitspray}.

It remains to prove the last claim in the lemma.
(This proof was omitted in \cite[Proposition 5.8.1]{Forstneric2017E} 
on the grounds of being trivial. 
We use this opportunity to include a sketch of proof.)
Assuming that the map $\psi(x,w)$ is tangent to $\psi_0(x,w)=w$ 
to order $k\ge 1$ along $w=0$, we must show that the same holds 
for the maps $a_\psi$ and $b_\psi$ in \eqref{eq:apsi}.
For simplicity, we consider the case when $w\in \C$ is a scalar variable;
a similar argument applies in general and we leave it to the reader.
We consider the Taylor expansions of our functions along $w=0$,
beginning with the case $k=1$. Thus, 
\[
	\psi(x,w)=h_1(x) w + O(|w|^2),\quad h_1\in \Lambda^r(C).
\]
Write 
\[
	c(x,w)          = \sum_{j=0}^\infty c_j(x)w^j, \quad 
	(\Acal c)(x,w) = \sum_{j=0}^\infty a_j(x)w^j, \quad
	(\Bcal c)(x,w) = \sum_{j=0}^\infty b_j(x)w^j, 
\]
where $c_j\in \Lambda^r_\Oscr(C)$, 
$a_j=\Acal c_j \in \Lambda^r_\Oscr(A)$, 
and $b_j=\Bcal c_j \in \Lambda^r_\Oscr(B)$ for all $j=0,1,\ldots$.
By the definition of $\Phi$ in \eqref{eq:Phi}, we have that 
\begin{eqnarray*}
	\Phi(\psi,c)(x,w) &=& 
	\psi\bigl(x,w+(\Acal c)(x,w)\bigr) - \bigl(w+(\Bcal c)(x,w)\bigr) \\
	&=&  h_1(x) \bigl(w+(\Acal c)(x,w)\bigr) - w - (\Bcal c)(x,w) + O(|w|^2)\\
	&=&  h_1(x) a_0(x) - b_0(x) + O(|w|) =0. 
\end{eqnarray*}
The equation $\Phi(\psi,c)=0$ gives $h_1a_0 -b_0=0$.
From $\Acal c-\Bcal c=c$ we get $a_0-b_0=c_0$, and hence 
$a_0(1-h_1)=c_0$ on $C$. Note that 
$\|a_0\|_{\Lambda^r(A)} \le const \|c_0\|_{\Lambda^r(C)}$. If 
$\psi(x,w)$ is close to $w$ then $h_1$ is close to $1$, which 
makes $\|a_0(1-h_1)\|_{\Lambda^r(C)}$ smaller than 
$\|c_0\|_{\Lambda^r(C)}$. In view of $a_0(1-h_1)=c_0$  
this forces $c_0=0$, and hence also $a_0=0$ and $b_0=0$, 
thereby proving the claim for $k=1$.

Assume now that $k>1$ and 
\[
	\psi(x,w)=w+ h_k(x) w^k + O(|w|^{k+1}),\quad 
	h_k\in\Lambda^r_\Oscr(C,\C^N). 
\]
The equation $\Phi(\psi,c)=0$ gives 
\begin{eqnarray*}
	\Phi(\psi,c)(x,w) &=& 
	\psi\bigl(x,w+(\Acal c)(x,w)\bigr) - \bigl(w+(\Bcal c)(x,w)\bigr) \\
	&=&  w+(\Acal c)(x,w) + h_k(x)(w+(\Acal c)(x,w))^k  
		- w - (\Bcal c)(x,w) + O(|w|^{k+1}) \\
	&=&  c(x,w) + h_k(x)(w+(\Acal c)(x,w))^k  + O(|w|^{k+1}) =0. 
\end{eqnarray*}
By the proof for $k=1$ we have $|(\Acal c)(x,w)|=O(|w|)$.
Therefore, $(w+(\Acal c)(x,w))^k=O(|w|^k)$ and the above identity
implies $c(x,w)=O(|w|^k)$. Thus, $c_j=0$ for $j=1,\ldots, k-1$, 
which also implies $a_j=0$ and $b_j=0$ for $j=1,\ldots, k-1$. 
\end{proof}


%
%  VERTICAL TANGENT BUNDLE OF A FIBRE BUNDLE
%
Assume now that $D$ is a compact strongly 
pseudoconvex domain with smooth boundary in a Stein manifold $X$,
$Y$ is a complex manifold, and $h:Z\to D$ is a topological fibre bundle
with fibre $Y$ which is holomorphic on $\mathring D=D\setminus bD$.
For $x\in D$ write $Z_x=h^{-1}(x)\cong Y$. The set
\[
	VT_z(Z) = T_z Z_{\pi(z)},\quad z\in Z 
\]
is called the vertical tangent space to $Z$ at $z$, 
and the vector bundle $VT(Z)\to Z$ with fibres $VT_z(Z)$ 
is the vertical tangent bundle of $(Z,h)$. 
When $h$ is differentiable in the variable $x\in D$, we have
$
	VT_z(Z) = \ker dh_z, \ z\in Z.
$
If the bundle $h:Z\to D$ is of H\"older--Zygmund class $\Lambda^r_\Oscr(D)$ 
(see Subsect.\ \ref{ss:Lr-bundles}) then $VT(Z)$ is of local 
class $\Lambda^r_\Oscr(Z)$, 
and for every section $f:D\to Z$ of class $\Lambda^r_\Oscr(D)$ the pullback
bundle $f^*VT(Z)\to D$ is a vector bundle of class $\Lambda^r_\Oscr(D)$.

We recall the notion of (local) dominating sprays of sections;
see \cite[Definition 5.9.1, p.\ 239]{Forstneric2017E} 
for sprays over open domains and \cite[Sect.\ 8.10]{Forstneric2017E} 
for sprays over compact domains in a complex manifold. 

%
%  SPRAYS
%
\begin{definition}\label{def:fibre-spray}
Let $h:Z\to D$ be a topological fibre bundle which is holomorphic on $\mathring D$. 
A holomorphic spray of sections of $h:Z\to D$  
is a continuous map $f : D\times W\to Z$ which is 
holomorphic on $\mathring D\times W$, where $0\in W\subset \C^N$ 
is an open convex set, such that 
\begin{equation} \label{eq:spray}
        h(f(x,w)) = x\quad \text{for $x \in D$ and $w\in W$}.
\end{equation}
The section $f_0=f(\cdotp,0):D\to Z$ is called the core of $f$.
The spray $f$ is dominating on a subset $K \subset D$ 
if the vertical derivative of $f$ at $w=0$, given by
\[
    	\di_w|_{w=0} f(x,w) :  T_0 \C^N\cong \C^N  \lra  VT_{f(x,0)} Z,
\]
is surjective for all $x\in K$. We say that $f$ is dominating if this 
holds for all $x\in D$. 
\end{definition}

We shall consider fibre bundles $h:Z\to D$ 
and sprays $f:D\times W\to Z$ of classes $\Lambda^r_\Oscr(D)$ 
for $r>0$. The main ingredient in the proof of Theorem \ref{th:main}
is the following gluing lemma for such sprays. Its analogue
in spaces $\Ascr^r(D)$, $r\in\Z_+$, was first proved in 
\cite[Proposition 4.3]{DrinovecForstneric2007DMJ}. See also 
the more concise proofs in \cite[Proposition 2.4]{DrinovecForstneric2008FM} 
and \cite[Proposition 5.9.2]{Forstneric2017E}. 

%
%
%  THE MAIN HEFTUNGSLEMMA
%
%
\begin{lemma}[Gluing sprays of class $\Lambda^r_\Oscr$]
\label{lem:gluing}
Let $(A,B)$ be a Cartan pair in a Stein manifold
(see Def.\ \ref{def:Cartan-pair}) and set $C=A\cap B$, $D=A\cup B$.
Assume that $h:Z\to D$ is a fibre bundle of class $\Lambda^r_\Oscr(D)$.
Given a convex domain $0\in W_0\subset\C^N$ and a spray of sections 
$f:A\times W_0\to Z$ of class $\Lambda^r_\Oscr(A\times W_0)$ 
which is dominating on $C$, there is a ball $W\subset\C^N$ 
with $0\in W \subset W_0$ satisfying the following conditions. 
%
\begin{enumerate}[\rm (a)]
\item 
For every holomorphic spray of sections
$g : B \times W_0 \to Z$ of class $\Lambda^r_\Oscr(B \times W_0)$
which is sufficiently close to $f$ in $\Lambda^r(C \times W_0)$
there exists a spray of sections $f' : D \times W \to Z$
of class $\Lambda^r_\Oscr(D\times W)$, close to $f$ in 
$\Lambda^r(A\times W)$ (depending on the $\Lambda^r$
distance between $f$ and $g$ on $C\times W_0$), 
whose core $f'_0$ is homotopic to $f_0=f(\cdotp,0)$ on $A$
and is homotopic to $g_0=g(\cdotp,0)$ on $B$.
%
\item
If $f$ and $g$ agree to order $k\in \N$ along $C\times\{0\}$, 
then $f'$ can be chosen to agree to order $k$ with $f$ 
along $A\times\{0\}$ and with $g$ along $B \times\{0\}$.
\end{enumerate}
%
The analogous result holds for families of sprays depending
continuously on a parameter $p\in P$ in a compact Hausdorff space.
If in addition the sprays $f$ and $g$ agree over $C$ for values of $p$ 
in a closed subset $Q\subset P$ which is a strong 
neighbourhood deformation retract, 
then $f'$ can be chosen to agree with $f$ and $g$ 
on $A$ and $B$, respectively, for $p\in Q$.
\end{lemma}

\begin{proof}[Sketch of proof]
The proof is analogous to that of 
\cite[Proposition 5.9.2, p.\ 240]{Forstneric2017E},
which applies to function spaces $\Ascr(D)$.
As noted in \cite[Remark 5.8.3 (B), p.\ 238]{Forstneric2017E},
the same proof applies in any Banach spaces on which there is a
linear bounded solution operator for the $\dibar$-equation 
on the level of $(0,1)$-forms. By Theorem \ref{th:dibar}, 
this holds for the Zygmund spaces $\Lambda^r$ with $r>0$. 
Nevertheless, some steps must be adjusted using results in  
Sections \ref{sec:approximation} and \ref{sec:vectorbundles}.
We indicate the necessary changes. 

The first step is to find a ball $0\in W\Subset W_0$ and a map 
$\gamma:C\times W\to C\times \C^N$ of the form 
\[ 
	\gamma(x,w) = (x,w + c(x,w)), \quad x\in C,\ w\in W,
\] 
with $c\in \Lambda^r_\Oscr(C\times W)$ close to $0$
(depending on the $\Lambda^r(C\times W_0)$ distance 
between $f$ and $g$), such that 
\begin{equation}\label{eq:fggama} 
	 f = g\circ \gamma \quad \text{holds on $C \times W$}.
\end{equation}
Such a map $\gamma$ is called a transition map between the sprays $f$ and $g$. 
(For maps of class $\Ascr^r$ with $r\in \Z_+$, 
a solution is given by \cite[Lemma 5.9.3, p.\ 240]{Forstneric2017E}.)  
In the process of constructing $\gamma$, 
we must split the trivial bundle $C\times \C^N$
by the subbundle $E'=\ker \di_w f|_{w=0}$,
the kernel of the vertical derivative of the spray $f$ at $w=0$. 
In our case, $E'$ is of class $\Lambda^r_\Oscr(C)$, 
and a splitting is furnished by Theorem \ref{th:complement} in the 
basic case and by Theorem \ref{th:complement-par} in the parametric case. 
The proof of \cite[Lemma 5.9.3]{Forstneric2017E} 
then applies without further changes. 

In the second step, pick a number $0<\epsilon<1$
and apply Lemma \ref{lem:splitting} to split $\gamma$ in the form 
\[
   \gamma\circ\alpha = \beta
   \ \ \text{on $C\times \epsilon W$},
\]
with $\alpha$ and $\beta$ as in \eqref{eq:alphabeta}. 
From this and \eqref{eq:fggama} it follows that 
\begin{equation}\label{eq:falphagbeta}
	f\circ \alpha = g\circ \beta \ \ \text{on $C\times \epsilon W$}.
\end{equation}
Hence, the two sides of the above equation amalgamate 
to a spray $f':D\times \epsilon W\to Z|_D$ satisfying the lemma. 
The proof in the parametric case follows the same scheme.
\end{proof}


%
%
%   PROOF OF THE MAIN THEOREM
%
%
\section{Proof of Theorem \ref{th:main}}\label{sec:proofmain}

We begin by explaining the proof of the basic
(nonparametric) case of Theorem \ref{th:main}.
The parametric version in Theorem \ref{th:mainpar}
includes the 1-parametric case, stated in Theorem \ref{th:main}. 

Thus, our task is to prove that any continuous section 
$f_0\in \Gamma(\overline \Omega,Z)$ of the fibre bundle 
$h:Z\to \overline \Omega$ of class $\Lambda^r_\Oscr(\overline \Omega)$ 
with an Oka fibre is homotopic to a section 
$f\in \Gamma^r_\Oscr(\overline \Omega,Z)$ which is holomorphic
on $\Omega$.

Choose a smooth strongly plurisubharmonic
function $\rho$ on a neighbourhood $U\subset X$ of $\overline\Omega$
such that 
\[
	\Omega=\{x\in U:\rho(x)<0\},\quad  
	d\rho_x\ne 0 \ \ \text{for every $x\in b\Omega =\{\rho=0\}$.}
\]
Pick $c<0$ such that $\rho$ has no critical values on the interval $[c,0]$.
Set 
\[
	A_0=\{\rho\le c\},\qquad 
	A'=\{\rho\le 0\}=\overline \Omega.
\]
Given an open cover $\Ucal=\{U_j\}$ of $\overline{A'\setminus A}$ 
consisting of holomorphic coordinate charts $U_j\subset X$,  
\cite[Lemma 5.10.3]{Forstneric2017E} gives compact, smoothly bounded, 
strongly pseudoconvex domains 
\begin{equation}\label{eq:sequenceAbis}
	A_0 \subset A_1 \subset \cdots \subset A_m=A'=\overline \Omega
\end{equation}
for some $m\in\N$ such that for every $k=0,1,\ldots,m-1$ we have 
$A_{k+1}=A_k\cup B_k$, where $(A_k,B_k)$ is a special Cartan pair
(see Definition \ref{def:Cartan-pair}) and $B_k\subset U_j$ for some $j=j(k)$.  
Hence, we may assume that the bundle $Z\to \overline \Omega$
is trivial over $B_k$ for all $k=0,1,\ldots,m-1$.

By the Oka principle on open Stein manifolds 
\cite[Theorem 5.4.4]{Forstneric2017E}, we may assume that 
$f_0$ is holomorphic on a neighbourhood of $A_0$. 
It suffices to show that for every special 
Cartan pair $(A,B)$ in $\overline \Omega$ such that the bundle
$Z\to \overline \Omega$ is trivial over $B$, 
we can approximate a section $f\in \Gamma^r_{\Oscr}(A,Z)$ as closely 
as desired by sections $\tilde f\in \Gamma^r_{\Oscr}(D,Z)$, 
where $D=A\cup B$. The theorem then follows by a finite induction,
using the sequence \eqref{eq:sequenceAbis} and starting with 
the section $f_0$ on $A_0$. The existence of a homotopy from $f_0$ 
to $f=f_m$ is obvious since there is no change of topology from $A_0$
to $A'=\overline \Omega$.

Fix a special Cartan pair $(A,B)$ and a section $f\in \Gamma^r_{\Oscr}(A,Z)$. 
The construction of a section $\tilde f \in \Gamma^r_{\Oscr}(D,Z)$ 
approximating $f$ in $\Gamma^r_{\Oscr}(A,Z)$ proceeds in three steps.
%
\begin{enumerate}[\rm (1)]
\item We embed $f$ as the core of a dominating spray of sections
$F:A\times W\to Z$ of class $\Lambda^r_{\Oscr}(A\times W,Z)$
(see Definition \ref{def:fibre-spray}), where $0\in W\subset \C^N$ 
is a ball. The construction of $F$ is explained in the sequel. 
%
\item 
Set $C=A\cap B$ and fix a number $\delta\in (0,1)$. 
Since the sets $C\subset B$ are convex in a local holomorphic coordinate
on a neighbourhood of $B$, the bundle $h$ is trivial over $B$, 
and the fibre $Y$ of $h$ is an Oka manifold, 
we can use Theorem \ref{th:approximation} to approximate $F$ 
as closely as desired in $\Lambda^r(C\times \delta W)$
by holomorphic sprays of sections $G:B\times \delta W \to Z$.
%
\item 
Assuming that the approximation in the previous step 
is close enough, we apply Lemma \ref{lem:gluing} to glue $F$ 
and $G$ into a spray $\wt F\in \Gamma^r_\Oscr(D\times \delta' W)$ 
for some $\delta'\in (0,\delta)$ which approximates $F$ 
in $\Lambda^r(A\times \delta' W)$. The section
$\tilde f:= \wt F(\cdotp,0) \in \Gamma^r_{\Oscr}(D,Z)$ 
then approximates $f$ in $\Gamma^r_{\Oscr}(A,Z)$. 
\end{enumerate}

This completes the proof, modulo the construction of a 
dominating spray in step (1). For this, we follow 
\cite[proof of Proposition 8.10.2, p.\ 388]{Forstneric2017E}
(the original reference is \cite[Proposition 4.1]{DrinovecForstneric2008FM}), 
where a result of this kind was proved for classes $\Ascr^r$ with $r\in\Z_+$.
Here is a sketch.

By Remark \ref{rem:TheoremA}, there is a vector bundle epimorphism 
\[
	L:A\times \C^N\longrightarrow f_0^* VT(Z)
\] 
of class $\Lambda^r_\Oscr(A\times\C^N)$ for some $N\in \N$.

\noindent {\bf Claim:} Given $L$ as above, there is a dominating spray 
$F:A\times W\to Z$ of class $\Lambda^r_{\Oscr}(A\times W,Z)$,
where $0\in W\subset \C^N$ is a ball, such that 
$F(\cdotp,0)=f_0$ and 
\begin{equation}\label{eq:L}
    	\di_w|_{w=0} F(x,w) = L(x,\cdotp): \C^N  \to  VT_{f(x,0)} Z
	\ \ \text{for all $x\in A$}.
\end{equation}

To see this, choose a sequence of compact, smoothly bounded, 
strongly pseudoconvex domains 
\[ 
	A_0 \subset A_1 \subset \cdots \subset A_m=A
\]
as in \eqref{eq:sequenceAbis} such that $A_0$ is contained in the interior
of $A$ and for every $k=0,1,\ldots,m-1$ we have 
$A_{k+1}=A_k\cup B_k$, where $(A_k,B_k)$ is a special Cartan pair
and the bundle $Z$ is trivial over $B_k$.  
The existence of a holomorphic spray $F_0:A_0\times W_0\to Z|_{A_0}$
with the core $F_0(\cdotp,0)=f_0|_{A_0}$ and satisfying \eqref{eq:L} 
is standard; see \cite[Proposition 8.10.2]{Forstneric2017E}. 
We inductively find sprays $F_k:A_k\times W_k \to Z|_{A_k}$ 
of classes $\Lambda^r_{\Oscr}(A_k\times W_k,Z)$ with the
core $f_0$ and satisfying \eqref{eq:L} for $x\in A_k$ $(k=1,\ldots,m)$, 
where $W_0\supset W_1\supset\cdots\supset W_m$ 
are balls in $\C^N$ centred at $0$. Every induction 
step is of the same kind and proceeds as follows.

We first approximate $F_k$ over $C_k=A_k\cap B_k$ 
by a spray $G_k:B_k\times V_k \to Z|_{B_k}$ of class 
$\Lambda^r_{\Oscr}(B_k\times V_k,Z)$ satisfying
\eqref{eq:L} for points $x\in B_k$, where $0\in V_k\subset W_k$ 
is a smaller ball; see \cite[Lemma 8.10.3]{Forstneric2017E} and note
that its proof also applies in classes $\Lambda^r$. 
In particular, $F_k$ and $G_k$ agree along  
$C\times \{0\}$ to the second order. Next, we apply 
Lemma \ref{lem:gluing} to glue $F_k$ and $G_k$ 
into a spray $F_{k+1}$ over $A_{k+1}=A_k\cup B_k$
with the core $f_0$ and satisfying condition \eqref{eq:L} for all 
points $x\in A_{k+1}$.
This is done by first finding a map $\gamma_k$
of the form \eqref{eq:gamma} and of class 
$\Lambda^r_\Oscr(C_k\times V_k)$ which agrees with the identity
to the second order along $C_k\times\{0\}$ and satisfies 
$F_k\circ \gamma_k = G_k$ on $C_k\times V_k$; see \eqref{eq:fggama}.
(The ball $V_k$ is allowed to shrink. The construction of such a map 
$\gamma_k$ was explained in the proof of Lemma \ref{lem:gluing}.) 
Next, we apply Lemma \ref{lem:splitting} to 
find sprays $\alpha_k$ and $\beta_k$ over $A_k$ and $B_k$
such that $\alpha_k \circ \gamma_k =\beta_k$ 
(see \eqref{eq:splitspray}),
and $\alpha_k$ and $\beta_k$ agree with the identity to the second
order along $A_k\times \{0\}$ and $B_k\times \{0\}$,
respectively. As in \eqref{eq:falphagbeta}, this yields
the spray $F_{k+1}$ over $A_{k+1}=A_k\cup B_k$, 
defined by the identity
\[
	F_k \circ \alpha_k  = G_k \circ \beta_k  \ \ 
	\text{on $C_k\times \epsilon V_k$},
	\quad 0<\epsilon<1.
\]
By the construction, $F_{k+1}$ has the core $f_0$
and satisfies condition \eqref{eq:L} on the set $A_{k+1}$. 
Taking $W_{k+1}=\epsilon V_k$ completes the induction step 
and proves the Claim. 

This completes the proof of the basic case of Theorem \ref{th:main}.

%
%  PARAMETRIC MAIN THEOREM
%
We now state a general parametric version of Theorem \ref{th:main}.
Recall that $\Gamma(\overline \Omega,Z)$ denotes the space of
continuous sections of a topological fibre bundle $h:Z\to \overline \Omega$. 
If the bundle is of class $\Lambda^r_\Oscr(\overline \Omega)$ then 
$\Gamma^r_\Oscr(\overline \Omega,Z)$ denotes the space of
sections of the same class.

\begin{theorem}\label{th:mainpar}
Assume that $\Omega$ is as in Theorem \ref{th:main},
$r>0$, $h:Z\to \overline \Omega$ is a fibre bundle of class 
$\Lambda^r_\Oscr(\overline \Omega)$ with Oka fibre, 
$P$ is a compact Hausdorff space, and $Q$ is a closed subset 
of $P$ which is a strong neighbourhood deformation retract.
Let $f:P\to \Gamma(\overline \Omega,Z)$ be a continuous map
such that $f|_Q:Q\to \Gamma^r_\Oscr(\overline \Omega,Z)$.
Then there is a homotopy 
$f_s:P\to \Gamma(\overline \Omega,Z)$, $s\in [0,1]$,
which is fixed on $Q$ such that $f_0=f$ and 
$f_1:P\to  \Gamma^r_\Oscr(\overline \Omega,Z)$.
\end{theorem}

The proof follows the same scheme as that of Theorem \ref{th:main},  
using the parametric versions of the tools 
in Sections \ref{sec:vectorbundles} and \ref{sec:gluing}.
(See also \cite[Sect.\ 5.13]{Forstneric2017E}.) 
We leave the details to the reader.

%
%
%   THE OKA PRINCIPLE FOR VECTOR BUNDLES 
%
%
\section{The Oka principle for vector bundles and principal
bundles of class $\Lambda^r_\Oscr$}
\label{sec:OPVB}

In this section, we prove Theorem \ref{th:OPVB}.
We follow \cite[proof of Theorem 5.3.1]{Forstneric2017E},
which is due to Grauert \cite{Grauert1958MA}. 
See also Cartan's exposition of Grauert's Oka principle in \cite{Cartan1958}.

\begin{proof}[Proof of Theorem \ref{th:OPVB}]
Denote by $G(m,N)$ the Grassmann manifold 
of all complex $m$-dimensional subspaces of $\C^N$.
A topological vector bundle $E \to \overline\Omega$ of rank $m$ is 
the pullback $f^*\U$ of the universal bundle $\U\to G(m,N)$ of rank $m$
by a continuous map $f:\overline\Omega\to G(m,N)$ for a suitable $N$.
(We take $N$ big enough such that $E$ embeds as a topological
vector subbundle of the trivial bundle $\overline\Omega \times \C^N$.)
Since $G(m,N)$ is a complex homogeneous manifold, and hence an 
Oka manifold by Grauert \cite{Grauert1957II}, 
Theorem \ref{th:main} shows that $f$ is 
homotopic to a map $F:\overline\Omega \to G(m,N)$ 
of class $\Lambda^r_\Oscr(\overline\Omega)$.
The pullback $F^*\U\to \overline\Omega$ is then a 
vector bundle of class $\Lambda^r_\Oscr(\overline\Omega)$
which is topologically isomorphic to $E\cong f^*\U$. 
This proves part (i) of the theorem.

To prove the second statement, let $E\to \overline\Omega$ and 
$E' \to \overline\Omega$ be vector bundles of class
$\Lambda^r_\Oscr(\overline\Omega)$ and rank $m$. 
There are an open cover $\{U_j\}$ of ${\overline\Omega}$ 
by smoothly bounded domains and vector bundle isomorphisms 
\[
	\theta_j: E|_{\overline U_j} \stackrel{\cong}{\lra} \overline U_j\times\C^m,
	\qquad
	\theta'_j: E'|_{\overline U_j} \stackrel{\cong}{\lra} 
	\overline U_j\times\C^m
\]
of class $\Lambda^r_\Oscr(\overline U_j)$. Set $U_{i,j}=U_i\cap U_j$. Let
\[
    g_{i,j} :  \overline U_{i,j} \to GL_m(\C), \qquad 
    g'_{i,j} : \overline U_{i,j} \to GL_m(\C)
\]
denote the fibrewise holomorphic transition maps of 
class $\Lambda^r_\Oscr(\overline U_{i,j})$ so that 
\[
    \theta_i \circ\theta_j^{-1}(x,v)= \bigl(x,g_{i,j}(x)v\bigr),
    \quad x \in \overline U_{i,j},\ v\in\C^m,
\]
and likewise for $E'$. A complex vector bundle isomorphism 
$\Phi : E \to E'$ is given by a collection of complex vector bundle 
isomorphisms 
\[ 
	\Phi_j: \overline U_j\times\C^m\to \overline U_j\times \C^m
\] 
of the form
\[
    \Phi_j(x,v)=\bigl(x,\phi_j(x)v\bigr),
    \quad x\in \overline U_j, \ v\in \C^m, 
\]
with $\phi_j(x)\in GL_m(\C)$ for $x\in \overline U_j$, 
satisfying the compatibility conditions
\begin{equation}\label{eq:bundle-P}
    \phi_i   = g'_{i,j}\phi_j g_{i,j}^{-1} 
    		= g'_{i,j} \phi_j g_{j,i} \quad {\rm on}\ \overline U_{i,j}.
\end{equation}
Let $h:Z\to {\overline\Omega}$ denote the fibre bundle of
class $\Lambda^r_\Oscr(\overline\Omega)$
with fibre $G=GL_m(\C)$ and transition maps \eqref{eq:bundle-P}. 
This means that $Z|_{\overline U_j}\cong \overline U_j\times G$ 
for each $j$, an element $(x,v) \in \overline U_j\times G$ 
for $x\in \overline U_{i,j}$ is identified with
$(x,v')\in \overline U_i\times G$ where $v'= g'_{i,j}(x)\, v \, g_{j,i}(x)$,
and no other identifications are made. 
A collection of maps $\phi_j : \overline  U_j\to G$ satisfying
conditions \eqref{eq:bundle-P} is then a section ${\overline\Omega}\to Z$. 
This shows that complex vector bundle isomorphisms $E\to E'$ 
correspond to sections of $Z\to {\overline\Omega}$, with 
isomorphisms of class $\Lambda^r_\Oscr(\overline\Omega)$
corresponding to sections of the same class. 
Hence, part (ii) follows from Theorem \ref{th:main}.
\end{proof}


Similarly one can prove the following analogue of 
\cite[Theorem 8.2.1]{Forstneric2017E} due to Grauert \cite{Grauert1958MA}.
In the proof, one should use the analogue of Lemma \ref{lem:Cartan} for an
arbitrary complex Lie group $G$; see Remark \ref{rem:Liegroup}.


\begin{theorem}\label{th:OPprincipalbundles}
Let $\Omega$ be as in Theorem \ref{th:OPVB}. 
For every complex Lie group $G$, the isomorphism classes of 
principal $G$ bundles on $\overline\Omega$ 
of class $\Lambda^r_\Oscr(\overline\Omega)$
are in bijective correspondence with the topological 
isomorphism classes of principal $G$ bundles.
\end{theorem}

%
%  GRASSMANN MANIFOLDS AND LU QI-KENG
% 
\begin{remark}\label{rem:Grassmann}
Grassmann manifolds $G(m,N)$ play a major role in the theory of 
complex vector bundles as the classifying spaces. Being projective, 
they are K\"ahler manifolds. An explicit formula
for a K\"ahler metric on $G(m,N)$ was given by Lu Qi-Keng in 
1963, see \cite{Lu-Qi-Keng1963I,Lu-Qi-Keng1963II}.
\end{remark}

% 
% 
% ACKNOWLEDGEMENTS 
% 
% 
%\noindent {\bf Statements and Declarations.}

\medskip \noindent
{\bf Acknowledgements.} 
Research was supported by the European Union 
(ERC Advanced grant HPDR, 101053085) and 
grants P1-0291 and N1-0237 from ARIS, Republic of Slovenia. 

I wish to thank Andrei Teleman for asking the question
which led to the main results of the paper (private communication,
January 2026), and for helpful communication regarding
the H\"older--Zygmund spaces. I also thank Xianghong Gong
for sharing with me his expertise on H\"older--Zygmund spaces.
A part of the background work on the paper originates in my 
joint papers \cite{DrinovecForstneric2007DMJ,DrinovecForstneric2008FM}
with Drinovec Drnov\v sek, whom I thank for this indirect contribution.
A part of the work was done during my visit to 
the University of Adelaide in February 2026. I wish to thank 
Finnur L\'arusson for the invitation and the mentioned 
institution for its hospitality.  

%
%  DECLARATIONS (IF NEEDED)
%
\begin{comment}
\noindent
{\bf Author Contributions} 
The corresponding author is the sole author of the paper, 
so the work is entirely accountable to him. 

\noindent
{\bf Data Availability} 
No datasets were generated or analysed during the current study. 

%\noindent{\bf Declarations}

\noindent
{\bf Conflicts of interest declaration}
The authors declare no conflicts of interest.

\noindent 
{\bf Competing interests declaration} The authors declare no competing interests.

\noindent 
{\bf Ethics, Consent to Participate, and Consent to Publish declarations:} not applicable.
\end{comment}



%%%%%%%%%% 
%%%%%%%%%% 
%%%%%%%%%% 
%%%%%%%%%% THE BIBLIOGRAPHY 
%%%%%%%%%% 
%%%%%%%%%% 

%{\bibliographystyle{abbrv} \bibliography{references}} 

%\begin{comment}

\begin{thebibliography}{10}

\bibitem{BealsGreinerStanton1987}
R.~Beals, P.~C. Greiner, and N.~K. Stanton.
\newblock {{\(L^ p\)}} and {Lipschitz} estimates for the
  {{\({\bar\partial}\)}}-equation and the {{\({\bar\partial}\)}}-{Neumann}
  problem.
\newblock {\em Math. Ann.}, 277:185--196, 1987.

\bibitem{BesovIlinNikolskii1979}
O.~V. Besov, V.~P. Il'in, and S.~M. Nikol'skii.
\newblock Integral representations of functions and imbedding theorems. {Vol}.
  {II}. {Ed}. by {Mitchell} {H}. {Taiblesson}. {Translation} from the
  {Russian}.
\newblock Scripta {Series} in {Mathematics}. {Washington}, {D}.{C}.: {V}.{H}.
  {Winston} \& {Sons}. {New} {York} etc.: {John} {Wiley} \& {Sons} {VIII}, 311
  p., 1979. 
  
\bibitem{BourdaudCristoforis2002}
G.~Bourdaud and M.~Lanza~de Cristoforis.
\newblock Functional calculus in {H{\"o}lder}-{Zygmund} spaces.
\newblock {\em Trans. Am. Math. Soc.}, 354(10):4109--4129, 2002.

\bibitem{Cartan1940}
H.~Cartan.
\newblock Sur les matrices holomorphes de {$n$} variables complexes.
\newblock {\em J. Math. Pures Appl.}, 19:1--26, 1940.

\bibitem{Cartan1958}
H.~Cartan.
\newblock Espaces fibr\'es analytiques.
\newblock In {\em Symposium internacional de topolog\'\i a algebraica
  ({I}nternational symposium on algebraic topology)}, pages 97--121.
  Universidad Nacional Aut\'onoma de M\'exico and UNESCO, Mexico City, 1958.

\bibitem{Catlin1988}
D.~Catlin.
\newblock A {Newlander}-{Nirenberg} theorem for manifolds with boundary.
\newblock {\em Mich. Math. J.}, 35(2):233--240, 1988.

\bibitem{ChenShaw2001}
S.-C. {Chen} and M.-C. {Shaw}.
\newblock {\em {Partial differential equations in several complex variables.}}
\newblock Providence, RI: Amer. Math. Soc.; Somerville, MA: International
  Press, 2001.

\bibitem{DrinovecForstneric2007DMJ}
B.~Drinovec~Drnov{\v{s}}ek and F.~Forstneri\v{c}.
\newblock Holomorphic curves in complex spaces.
\newblock {\em Duke Math. J.}, 139(2):203--253, 2007.

\bibitem{DrinovecForstneric2008FM}
B.~Drinovec~Drnov{\v{s}}ek and F.~Forstneri\v{c}.
\newblock Approximation of holomorphic mappings on strongly pseudoconvex
  domains.
\newblock {\em Forum Math.}, 20(5):817--840, 2008.

\bibitem{FollandKohn1972}
G.~B. Folland and J.~J. Kohn.
\newblock {\em The {Neumann} problem for the {Cauchy}-{Riemann} complex},
  volume~75 of {\em Ann. Math. Stud.}
\newblock Princeton University Press, Princeton, NJ, 1972.

\bibitem{FornaessForstnericWold2020}
J.~E. {Forn{\ae}ss}, F.~{Forstneri\v{c}}, and E.~{Wold}.
\newblock {Holomorphic approximation: the legacy of Weierstrass, Runge,
  Oka-Weil, and Mergelyan.}
\newblock In {\em {Advancements in complex analysis. From theory to practice}},
  pages 133--192. Cham: Springer, 2020.

\bibitem{Forstneric2023Indag}
F.~Forstneri{\v{c}}.
\newblock Recent developments on {Oka} manifolds.
\newblock {\em Indag. Math., New Ser.}, 34(2):367--417, 2023.

\bibitem{Forstneric2007AJM}
F.~Forstneri\v{c}.
\newblock Manifolds of holomorphic mappings from strongly pseudoconvex domains.
\newblock {\em Asian J. Math.}, 11(1):113--126, 2007.

\bibitem{Forstneric2009CR}
F.~Forstneri\v{c}.
\newblock Oka manifolds.
\newblock {\em C. R. Math. Acad. Sci. Paris}, 347(17-18):1017--1020, 2009.

\bibitem{Forstneric2017E}
F.~Forstneri\v{c}.
\newblock {\em Stein manifolds and holomorphic mappings (The homotopy principle
  in complex analysis)}, volume~56 of {\em Ergebnisse der Mathematik und ihrer
  Grenzgebiete. 3. Folge}.
\newblock Springer, Cham, second edition, 2017.

\bibitem{Forstneric2026ICM}
F.~Forstneri{\v c}.
\newblock From Stein manifolds to Oka manifolds: the h-principle in complex
  analysis.
\newblock To appear in Proc. ICM 2026. \url{https://arxiv.org/abs/2509.21197}.

\bibitem{ForstnericLarusson2011}
F.~Forstneri\v{c} and F.~L{\'a}russon.
\newblock Survey of {O}ka theory.
\newblock {\em New York J. Math.}, 17A:11--38, 2011.

\bibitem{ForstnericLarusson2025MRL}
F.~Forstneri\v{c} and F.~L{\'a}russon.
\newblock Every projective {O}ka manifold is elliptic.
\newblock {\em Math. Res. Lett.}, 33(1):297--312, 2026.

\bibitem{ForstnericLowOvrelid2001}
F.~Forstneri\v{c}, E.~L{\o}w, and N.~{\O}vrelid.
\newblock Solving the {$d$}- and {$\overline\partial$}-equations in thin tubes
  and applications to mappings.
\newblock {\em Michigan Math. J.}, 49(2):369--416, 2001.

\bibitem{GanGong2024}
C.~Gan and X.~Gong.
\newblock Global {Newlander}-{Nirenberg} theorem for domains with {{\(C^2\)}}
  boundary.
\newblock {\em Mich. Math. J.}, 74(2):283--329, 2024.

\bibitem{GilbargTrudinger1983}
D.~Gilbarg and N.~S. Trudinger.
\newblock {\em Elliptic partial differential equations of second order. 2nd
  ed}, volume 224 of {\em Grundlehren Math. Wiss.}
\newblock Springer, Cham, 1983.

\bibitem{GohbergLancasterRodman1986}
I.~Gohberg, P.~Lancaster, and L.~Rodman.
\newblock {\em Invariant subspaces of matrices with applications}.
\newblock Can. Math. Soc. Ser. Monogr. Adv. Texts. John Wiley, New York, NY,
  1986.

\bibitem{Gong2025TAMS}
X.~Gong.
\newblock On regularity of $\overline\partial$-solutions on $a_q$ domains with
  $c^2$ boundary in complex manifolds.
\newblock {\em Trans. Amer. Math. Soc.}, 378(3):1771--1829.

\bibitem{Gong2019MA}
X.~Gong.
\newblock H{\"o}lder estimates for homotopy operators on strictly pseudoconvex
  domains with {{\(C^2\)}} boundary.
\newblock {\em Math. Ann.}, 374(1-2):841--880, 2019.

\bibitem{GongShi2024}
X.~Gong and Z.~Shi.
\newblock Global {N}ewlander--{N}irenberg theorem on domains with finite smooth
  boundary in complex manifolds.
\newblock {\em Sci.\ China Math.}, 69, 2026.
\newblock \url{https://doi.org/10.1007/s11425-025-2584-7}  
% https://www.sciengine.com/SCM/doi/10.1007/s11425-025-2584-7   

\bibitem{Grauert1957I}
H.~Grauert.
\newblock Approximationss\"atze f\"ur holomorphe {F}unktionen mit {W}erten in
  komplexen {R}\"aumen.
\newblock {\em Math. Ann.}, 133:139--159, 1957.

\bibitem{Grauert1957II}
H.~Grauert.
\newblock Holomorphe {F}unktionen mit {W}erten in komplexen {L}ieschen
  {G}ruppen.
\newblock {\em Math. Ann.}, 133:450--472, 1957.

\bibitem{Grauert1958MA}
H.~Grauert.
\newblock Analytische {F}aserungen \"uber holomorph-vollst{\"a}ndigen
  {R}{\"a}umen.
\newblock {\em Math. Ann.}, 135:263--273, 1958.

\bibitem{Gromov1989}
M.~Gromov.
\newblock Oka's principle for holomorphic sections of elliptic bundles.
\newblock {\em J. Amer. Math. Soc.}, 2(4):851--897, 1989.

\bibitem{GunningRossi1965}
R.~C. Gunning and H.~Rossi.
\newblock {\em Analytic functions of several complex variables}.
\newblock 1965.

\bibitem{HenkinLeiterer1984}
G.~M. Henkin and J.~Leiterer.
\newblock {\em Theory of functions on complex manifolds}, volume~60 of {\em
  Mathematische Lehrb\"ucher und Monographien, II. Abteilung: Mathematische
  Monographien [Mathematical Textbooks and Monographs, Part II: Mathematical
  Monographs]}.
\newblock Akademie-Verlag, Berlin, 1984.

\bibitem{Heunemann1986I}
D.~Heunemann.
\newblock An approximation theorem and {O}ka's principle for holomorphic vector
  bundles which are continuous on the boundary of strictly pseudoconvex
  domains.
\newblock {\em Math. Nachr.}, 127:275--280, 1986.

\bibitem{Heunemann1986III}
D.~Heunemann.
\newblock Extension of the complex structure from {S}tein manifolds with
  strictly pseudoconvex boundary.
\newblock {\em Math. Nachr.}, 128:57--64, 1986.

\bibitem{Heunemann1986II}
D.~Heunemann.
\newblock Theorem {B} for {S}tein manifolds with strictly pseudoconvex
  boundary.
\newblock {\em Math. Nachr.}, 128:87--101, 1986.

\bibitem{Kohn1963}
J.~J. Kohn.
\newblock Harmonic integrals on strongly pseudo-convex manifolds. {I}.
\newblock {\em Ann. of Math. (2)}, 78:112--148, 1963.

\bibitem{Kohn1964}
J.~J. Kohn.
\newblock Harmonic integrals on strongly pseudo-convex manifolds. {II}.
\newblock {\em Ann. of Math. (2)}, 79:450--472, 1964.

\bibitem{Kusakabe2021IUMJ}
Y.~Kusakabe.
\newblock Elliptic characterization and localization of {O}ka manifolds.
\newblock {\em Indiana Univ. Math. J.}, 70(3):1039--1054, 2021.

\bibitem{Kusakabe2024AM}
Y.~Kusakabe.
\newblock Oka properties of complements of holomorphically convex sets.
\newblock {\em Ann. Math. (2)}, 199(2):899--917, 2024.

\bibitem{Leiterer1976I}
J.~Leiterer.
\newblock Analytische {F}aserb\"undel mit stetigem {R}and {\"u}ber streng
  pseudokonvexen {G}ebieten. {I}. {G}arbentheoretische {H}ilfsmittel.
\newblock {\em Math. Nachr.}, 71:329--344, 1976.

\bibitem{Leiterer1976II}
J.~Leiterer.
\newblock Analytische {F}aserb\"undel mit stetigem {R}and \"uber streng
  pseudokonvexen {G}ebieten. {II}. {T}opologische {K}lassifizierung.
\newblock {\em Math. Nachr.}, 72:201--217, 1976.

\bibitem{Leiterer1990}
J.~{Leiterer}.
\newblock {Holomorphic vector bundles and the Oka-Grauert principle}.
\newblock In {\em Several complex variables. IV. Algebraic aspects of complex
  analysis}, Encycl. Math. Sci., Vol. 10, pages 63--103. Springer-Verlag, 1990.
\newblock Translation from {\em Itogi Nauki Tekh., Ser. Sovrem. Probl. Mat.,
  Fundam. Napravleniya} 10, 75-121 (1986). Translated by D. N. Akhiezer.

\bibitem{Lu-Qi-Keng1963I}
Q.-K. Lu.
\newblock The elliptic geometry of extended spaces.
\newblock {\em Acta Math. Sin.}, 13:49--62, 1963.

\bibitem{Lu-Qi-Keng1963II}
Q.-K. Lu.
\newblock A note about '{The} elliptic geometry of extended spaces'.
\newblock {\em Acta Math. Sin.}, 13:314, 1963.

\bibitem{Malgrange1966}
B.~Malgrange.
\newblock {\em Ideals of differentiable functions}, volume~3 of {\em Tata Inst.
  Fundam. Res., Stud. Math.}
\newblock London: Oxford University Press, 1966.

\bibitem{McShane1934}
E.~J. McShane.
\newblock Extension of range of functions.
\newblock {\em Bull. Am. Math. Soc.}, 40:837--842, 1934.

\bibitem{Ohsawa1984LNM}
T.~Ohsawa.
\newblock Holomorphic embedding of compact s.p.c.\ manifolds into complex
  manifolds as real hypersurfaces.
\newblock In {\em Differential geometry of submanifolds ({K}yoto, 1984)},
  volume 1090 of {\em Lecture Notes in Math.}, pages 64--76. Springer, Berlin,
  1984.

\bibitem{Oka1939}
K.~{Oka}.
\newblock {Sur les fonctions analytiques de plusieurs variables. III.
  Deuxi\`eme probl\`eme de Cousin.}
\newblock {\em {J. Sci. Hiroshima Univ., Ser. A}}, 9:7--19, 1939.

\bibitem{Palais1968}
R.~S. Palais.
\newblock {\em Foundations of global non-linear analysis}.
\newblock W. A. Benjamin, Inc., New York-Amsterdam, 1968.

\bibitem{Palais1971}
R.~S. Palais.
\newblock Banach manifolds of fiber bundle sections.
\newblock In {\em Actes du {C}ongr\`es {I}nternational des {M}ath\'ematiciens
  ({N}ice, 1970), {T}ome 2}, pages 243--249. Gauthier-Villars, Paris, 1971.

\bibitem{Rychkov1999}
V.~S. Rychkov.
\newblock On restrictions and extensions of the {Besov} and
  {Triebel}-{Lizorkin} spaces with respect to {Lipschitz} domains.
\newblock {\em J. Lond. Math. Soc., II. Ser.}, 60(1), 1999.

\bibitem{Seeley1964}
R.~Seeley.
\newblock Extension of {{\(C^ \infty\)}} functions defined in a half space.
\newblock {\em Proc. Am. Math. Soc.}, 15:625--626, 1964.

\bibitem{ShiYao2024}
Z.~Shi and L.~Yao.
\newblock New estimates of {Rychkov}'s universal extension operator for
  {Lipschitz} domains and some applications.
\newblock {\em Math. Nachr.}, 297(4):1407--1443, 2024.

\bibitem{ShiYao2025}
Z.~Shi and L.~Yao.
\newblock Sobolev {{\(\frac{1}{2}\)}} estimate for the {{\(\overline{\partial}\)}} 
equation on strictly pseudoconvex domains with {{\(C^2\)}} boundary.
\newblock {\em Am. J. Math.}, 147(1):235--277, 2025.

\bibitem{SteinEM1970}
E.~M. Stein.
\newblock {\em Singular integrals and differentiability properties of
  functions}, volume~30 of {\em Princeton Math. Ser.}
\newblock Princeton University Press, Princeton, NJ, 1970.

\bibitem{TaylorME2023}
M.~E. Taylor.
\newblock {\em Partial differential equations {I}. {Basic} theory}, volume 115
  of {\em Appl. Math. Sci.}
\newblock Cham: Springer, 3rd corrected and expanded edition edition, 2023.

\bibitem{Teleman2026}
A.~Teleman. 
\newblock An extension theorem for bundles with respect to strictly pseudoconvex extensions. 
\newblock  {\em arXiv e-print 2608.26404}, August 2026. 
\newblock \url{https://arxiv.org/abs/2608.26404}.

\bibitem{Wallin1988}
H.~Wallin.
\newblock New and old function spaces.
\newblock Function spaces and applications, {Proc}. {US}-{Swed}. {Semin}.,
  {Lund}/{Swed}., {Lect}. {Notes} {Math}. 1302, 99-114 (1988), 1988.

\bibitem{Whitney1934TAMS}
H.~Whitney.
\newblock Analytic extensions of differentiable functions defined in closed
  sets.
\newblock {\em Trans. Amer. Math. Soc.}, 36(1):63--89, 1934.

\end{thebibliography}
%\end{comment}

\end{document}

